% Standalone source generated from article.tex and the modular files.
% BEGIN INCLUDED SOURCE: preamble.tex
\documentclass[11pt,a4paper]{article}
\usepackage[T1]{fontenc}
\usepackage[utf8]{inputenc}
\usepackage{lmodern}
\usepackage[margin=24mm,headheight=15pt]{geometry}
\usepackage{amsmath,amssymb,amsthm,mathtools,mathrsfs}
\usepackage{microtype}
\usepackage{booktabs,longtable,array,enumitem}
\usepackage[dvipsnames]{xcolor}
\usepackage{fancyhdr}
\usepackage{hyperref}
\usepackage{bookmark}
\hypersetup{colorlinks=true,linkcolor=MidnightBlue,citecolor=MidnightBlue,
 urlcolor=MidnightBlue,pdftitle={Alternating Tails and Higher Spectral Defects},
 pdfauthor={Research continuation prepared for the ProveIt project},
 pdfsubject={Exact polylogarithm, Gamma, harmonic, zeta, Stieltjes and moment identities}}
\urlstyle{same}
\numberwithin{equation}{section}
\newtheorem{theorem}{Theorem}[section]
\newtheorem{lemma}[theorem]{Lemma}
\newtheorem{proposition}[theorem]{Proposition}
\newtheorem{corollary}[theorem]{Corollary}
\theoremstyle{definition}
\newtheorem{definition}[theorem]{Definition}
\newtheorem{question}{Question}
\theoremstyle{remark}
\newtheorem{remark}[theorem]{Remark}
\newcommand{\repo}{\texttt{VladimirReshetnikov/ProveIt}}
\newcommand{\sourcepin}{\texttt{73f18ef013d5b74eafc883efb2c15210642df12a}}
\newcommand{\shortpin}{\texttt{73f18ef013d5}}
\newcommand{\dd}{\,\mathrm d}
\setlength{\parskip}{3pt}
\setlength{\parindent}{1.25em}
\setlist{itemsep=4pt,topsep=5pt}
\allowdisplaybreaks[2]
\emergencystretch=1em
\pagestyle{fancy}
\fancyhf{}
\fancyhead[L]{\small Alternating tails and higher spectral defects}
\fancyhead[R]{\small ProveIt research continuation}
\fancyfoot[C]{\thepage}
\renewcommand{\headrulewidth}{0.3pt}


% END INCLUDED SOURCE: preamble.tex

\begin{document}
\hypersetup{pageanchor=false}
\begin{titlepage}
\vspace*{13mm}
\noindent{\small\sffamily\color{MidnightBlue} PROVEIT \quad / \quad RESEARCH CONTINUATION}
\vspace{12mm}

\noindent{\LARGE\bfseries Alternating Tails and\\[4pt] Higher Spectral Defects\par}
\vspace{7mm}
\noindent{\large Cyclotomic generators, all-index Stieltjes identities,\\[3pt]
Gamma primitives, and exact moment-reconstruction certificates\par}
\vspace{10mm}
\noindent Prepared for Vladimir Reshetnikov's ProveIt programme\\
Mathematical derivations and computational audit with OpenAI assistance\\[5pt]
11 October 2026
\vspace{11mm}

\begin{abstract}
We develop four exact continuations of the current polylogarithm
manuscript and its incoming reports. A finite adjacent-tail certificate
gives a complementary square identity; its continuous-order version
produces ordinary convergent identities for every generalized Stieltjes
index and an anchored Gamma primitive. A finite cyclotomic generator
reduces the prescribed conjugation sector of rational-tail products at
every conductor to trigonometric functions and digamma values, with
explicit conductor-four evaluations. For a quadratic spectrum with an
arbitrary polynomial weight, we compute every spectral derivative of
the conjugate-factor defect. The answer separates an ordinary
convergent logarithmic sum from finite Hurwitz and Stieltjes terms,
including the pole contribution that is lost by regular-part
differentiation. For any number of quadratic factors and every even
polynomial weight, we determine the complete second derivative, a
digamma double primitive of the mass correction, and all coincident-mass
coefficients. Finally, exact Bernoulli congruences prove new infinite
families of optimal moment prefixes at orders eight and ten, leaving
one explicitly identified residue class at each order unresolved.

The proofs are analytic or exact algebraic. Symbolic certificates and
independent high-precision diagnostics are supplied separately, with
their different evidential roles stated. The article reconciles source
status, preserves existing corrections, and proposes concrete further
research questions. Novelty is relative to the inspected project
corpus; global literature priority is not asserted.
\end{abstract}
\vfill
\noindent\textbf{Repository anchor.} \repo, commit \shortpin.\\
\noindent\textbf{Scope.} Exact identities and proofs; no new formalization claim.\\
\noindent\textbf{Companion artifacts.} Modular and standalone \LaTeX,
exact certificates, verification programs, recorded results, and integration notes.
\end{titlepage}
\pagenumbering{roman}
\hypersetup{pageanchor=true}
\begingroup
\small
\setlength{\parskip}{0pt}
\tableofcontents
\endgroup
\clearpage
\pagenumbering{arabic}
% BEGIN INCLUDED SOURCE: sections/01_scope.tex
\section{Scope, conventions, and the results established}
\label{intro:section}

This article continues the exact-identity programme in the unified
ProveIt manuscript and the six archives present in its incoming
directory at commit
\begin{center}\sourcepin.\end{center}
The commit is dated 11 October 2026, 18:22:01 UTC. The incoming archive
interiors were retrieved and inspected for overlap with the selected
targets. The repository manuscript and the relevant primary references
were checked where the proofs use them. This is a targeted mathematical
audit of those dependencies, not an exhaustive certification of every
statement in the repository. The deliverable is a separate article;
repository integration can preserve the existing thematic organization.

The main technical theme is to retain exact finite differences and
finite asymptotic terms before taking a limit. For tails, the missing
information is an endpoint product. For zeta continuation, it is a
finite residue contribution. For moment reconstruction, it is an exact
rational determinant rather than a floating rank test. These three
mechanisms lead to complementary identities among the functions named
in the research programme.

\subsection{A map from source questions to the present results}

\begin{longtable}{@{}>{\raggedright\arraybackslash}p{0.24\textwidth}>{\raggedright\arraybackslash}p{0.46\textwidth}>{\raggedright\arraybackslash}p{0.24\textwidth}@{}}
\toprule
Source target & Result established here & Precise status\\
\midrule
\endfirsthead
\toprule Source target & Result established here & Precise status\\\midrule
\endhead
Harmonic-tail square identities & A rational-tail finite certificate,
continuous order, all-index Stieltjes products, and anchored Gamma
integration (Section~\ref{ats:section}) & Additional identities with
ordinary convergence.\\[5pt]
Shifts and Twists, question Q6 \cite{HTST} & A finite generator for
arbitrary root-of-unity inner and outer colors
(Theorem~\ref{cyc:master}) & Solved for the specified parity/conjugation
sector; the complementary sector remains open.\\[5pt]
Quadratic-Spectrum Stieltjes, question 2 \cite{QSpec} & Every order
$\Delta_r$ for every polynomial weight
(Theorem~\ref{sd:thm:all-orders}) & Complete explicit correction formula;
generic finite polylogarithm reduction is not claimed.\\[5pt]
The same report, question 7 & Any number of quadratic factors, every
even polynomial weight, and every coincident-mass coefficient
(Theorems~\ref{mm:thm:second}--\ref{mm:thm:confluence}) & Complete at the
second spectral derivative in this setting.\\[5pt]
Optimal moment prefix, order eight \cite{PRR} & Full formal rank for
$p=8$, $q=2b$, $b\ge5$, outside $8001\pmod{16384}$
(Theorem~\ref{rt:rank8}) & Strengthens the existing exclusion
$1\pmod{16}$; one smaller progression remains.\\[5pt]
Optimal moment prefix, order ten & Full formal rank for $p=10$,
$q=2b$, $b\ge6$, outside $962\pmod{2048}$
(Theorem~\ref{rt:rank10}) & New infinite family and exact finite base
cases; no all-order theorem is inferred.\\
\bottomrule
\end{longtable}

The gap-four moment-prefix theorem is already established in the two
latest harmonic-tail reports \cite{HTGC,HTST}. It is not reproved here
as a new result. The Gaussian weight-six and revised weight-eight
numerical conjectures are not settled by this article. The new
cyclotomic identities concern different, explicitly specified sums;
their existence is not a missing proof of those numerical equalities.

\subsection{Conventions fixed throughout}

The Hurwitz zeta function is initially
$\zeta(s,a)=\sum_{n\ge0}(n+a)^{-s}$ for $\Re s>1$, $\Re a>0$,
then continued in its usual variables. Our generalized Stieltjes
constants are defined by
\begin{equation}\label{intro:stieltjes}
 \zeta(1+u,a)=\frac1u+
       \sum_{r\ge0}\frac{(-1)^r}{r!}\gamma_r(a)u^r.
\end{equation}
Thus $\gamma_0(a)=-\psi(a)$, $\gamma_0(1)=\gamma$, and a literal
linear Laurent coefficient is $-\gamma_1(a)$.
The prime in $\zeta'(s,a)$ always differentiates $s$; a derivative
in $a$ is written $\partial_a$. We use
\[
 \psi(a)=\frac{\Gamma'(a)}{\Gamma(a)},\qquad
 \psi^{(r)}(a)=(-1)^{r+1}r!\zeta(r+1,a)\quad(r\ge1),
\]
and $H_n^{(r)}=\sum_{k=1}^{n}k^{-r}$, with $H_n=H_n^{(1)}$ and
$H_0=0$. The polylogarithm is
$\operatorname{Li}_s(z)=\sum_{k\ge1}z^k/k^s$ in its convergence
disk, with continuation or ordinary boundary values stated when used.
These conventions agree with the standard references
\cite{DLMFZeta,DLMFPolylog,DLMFPolygamma}.

All complex powers use principal logarithms on the initial positive
shift domains. Any later continuation or antiderivative specifies its
domain and its anchor. A square of a complex tail means its algebraic
square; a modulus square is written explicitly. A series called
\emph{ordinary} is a limit of its stated partial sums. A continued
spectral value is never silently substituted for a divergent ordinary
sum. ``Locally normally convergent'' means locally uniform absolute
convergence, sufficient for holomorphic parameter differentiation.

The symbols used for colors in Section~\ref{cyc:sec} are roots of
unity; the positive real shift and mass variables in
Sections~\ref{sd:section}--\ref{mm:section} are local to those sections.
Nielsen's alternating Hurwitz function $\mathcal B_s$ is distinguished
from the Bernoulli polynomials $B_j$ and the cyclotomic generating
function $B_\omega$. The moment-rank statements concern a formal
coordinate space in which only the specified exact reduction relations
have been imposed. They make no arithmetic independence assertion
about actual zeta values.

\subsection{What counts as proof in the companion package}

The infinite-dimensional assertions below have written analytic proofs,
including convergence, pole cancellation, and branch control. Exact
symbolic certificates verify the finite algebra used in those proofs.
The moment nonvanishing theorems combine a uniform Bernoulli congruence
with exact polynomial identities and finitely many rational base cases.
High-precision evaluations supply independent diagnostics for signs,
factors, and endpoint conventions. They are not interval certificates,
proofs of an infinite family, or proofs of period independence.
No statement in this delivery is claimed to have been formalized in
Lean or another proof assistant.


% END INCLUDED SOURCE: sections/01_scope.tex

% BEGIN INCLUDED SOURCE: sections/02_tail_squares.tex
% Research component; no standalone preamble.  Prepared for the parent article.
% Sources: latest Shifts/Twists incoming report, section 05_twisted.tex;
% background: Hoffman, arXiv:1602.03399; DLMF 5.5 and 25.11;
% Berg--Koumandos--Pedersen, arXiv:1905.04131 (Nielsen beta notation).

\section{A complementary square law from one finite difference}
\label{ats:section}

The incoming alternating-outer resolvent theorem \cite{HTST} evaluates
$\sum_{n\ge0}(-1)^n(2n+1)E_n^\alpha(x)E_n^\beta(y)$ through
alternating double-zeta parity.  There is a complementary ordinary-outer
square law, with a direct finite certificate.  Its continuous-order
extension yields ordinary convergent identities involving generalized
Stieltjes constants.  The finite-difference mechanism is elementary;
the results below are additions relative to the inspected project
sources, not claims of literature priority.  Products of zeta tails are
a classical subject; see Hoffman's multiple-zeta treatment \cite{Hoffman}.

\subsection{The exact finite certificate and its compatible color}

Let $P$ be a polynomial of degree $d\ge2$ with $P(k)\ne0$ for every
positive integer $k$. Let $0<|\lambda|\le1$, $\lambda\ne1$, and define
\[
 E_n=\sum_{k>n}\frac{\lambda^k}{P(k)},\qquad
 t=-\lambda^{-1},\qquad
 Q_n=t^nP(n+1)E_nE_{n+1}\quad(n\ge0).
\]
The value $P(0)=0$ is allowed. All squares and products are algebraic;
complex conjugation is not inserted.

\begin{theorem}[A square identity for rational tails]
\label{ats:thm:polynomial}
The following series is absolutely convergent:
\begin{equation}\label{ats:eq:polynomial}
 \sum_{n\ge0}t^n\{P(n+1)+\lambda P(n)\}E_n^2
 =\lambda E_0\{1+P(0)E_0\}.
\end{equation}
More precisely, for every integer $N\ge0$,
\begin{equation}\label{ats:eq:finite}
 \sum_{n=0}^{N}t^n\{P(n+1)+\lambda P(n)\}E_n^2
 =\lambda E_0\{1+P(0)E_0\}+Q_N.
\end{equation}
It holds locally normally in the coefficients of $P$ on sets of fixed
degree avoiding the positive-integer poles, for each fixed admissible
$\lambda$.
\end{theorem}

\begin{proof}
For $n\ge1$, substitute
$E_{n-1}=E_n+\lambda^n/P(n)$ and
$E_{n+1}=E_n-\lambda^{n+1}/P(n+1)$ into the difference of two
successive $Q$'s. Since $t^{-1}=-\lambda$, the linear-tail terms cancel:
\[
 Q_n-Q_{n-1}
 =t^n\{P(n+1)+\lambda P(n)\}E_n^2.
\]
At the lower endpoint,
$Q_0=P(1)E_0^2-\lambda E_0$, whereas the $n=0$ summand is
$\{P(1)+\lambda P(0)\}E_0^2$. Their difference is exactly the
right side of \eqref{ats:eq:polynomial}. This proves the finite identity
without evaluating $1/P(0)$.

When $|\lambda|<1$, $E_n=O(|\lambda|^{n+1}n^{-d})$; the factor
$|t|^n$ therefore leaves geometric decay. On the unit circle with
$\lambda\ne1$, summation by parts and the bounded partial sums of
$\lambda^k$ give $E_n=O(n^{-d})$: use
$P(k)^{-1}=O(k^{-d})$ and
$P(k+1)^{-1}-P(k)^{-1}=O(k^{-d-1})$.
Thus the summands are $O(n^{-d})$, and $Q_N=O(N^{-d})\to0$.
These estimates are uniform on compact coefficient sets as stated.
\end{proof}

The color constraint is forced within this short certificate class.
If the same ansatz for $Q_n$ is made with arbitrary $t\ne0$, its
difference is
\[
 t^n\left[\{P(n+1)-t^{-1}P(n)\}E_n^2
      -\lambda^n(\lambda+t^{-1})E_n\right].
\]
Hence cancellation of the linear-tail term as an identity requires
$t=-\lambda^{-1}$. This is a classification of the displayed ansatz,
not an obstruction to other telescopers.

In particular, taking $P(k)=k^d$ gives an analytic polylogarithmic
identity for every integer $d\ge2$:
\begin{align}\label{ats:eq:polylog-tail}
 &\sum_{n\ge0}(-\lambda^{-1})^n\{(n+1)^d+\lambda n^d\}
   \left\{\operatorname{Li}_d(\lambda)
               -\sum_{k=1}^{n}\frac{\lambda^k}{k^d}\right\}^{\!2}
 =\lambda\operatorname{Li}_d(\lambda).
\end{align}
At $\lambda=0$ both sides have a removable value zero; the tails
contain $\lambda^{n+1}$, so no negative power survives in a summand.
For example, $d=2$ and $\lambda=i$ give the Gaussian identity
\[
 \sum_{n\ge0}i^n\{(n+1)^2+in^2\}
 \left\{\operatorname{Li}_2(i)-\sum_{k=1}^{n}\frac{i^k}{k^2}\right\}^{\!2}
 =-G-\frac{i\pi^2}{48},
\]
where $G=\beta(2)$ is Catalan's constant. This follows from
$\operatorname{Li}_2(i)=-\pi^2/48+iG$.

For $\lambda=-1$, $t=1$ and the weight is $P(n+1)-P(n)$, whose
degree is at most $d-1$. The same proof in this case permits $d=1$;
the summands are then $O(n^{-2})$.

At a root of unity of order $N$, the endpoint $E_0$ has a finite
polygamma reduction. Indeed, partial fractions reduce $1/P(k)$ to
terms $(k+z)^{-j}$, and residue classes give
\[
 \sum_{k=1}^{\infty}\frac{\lambda^k}{(k+z)^j}
 =N^{-j}\sum_{r=1}^{N}\lambda^r
                   \zeta\!\left(j,\frac{r+z}{N}\right)
 \qquad(j\ge2).
\]
For $j=1$ the correct pole-cancelled interpretation is
\[
 -\frac1N\sum_{r=1}^{N}\lambda^r
                         \psi\!\left(\frac{r+z}{N}\right),
\]
because $\sum_{r=1}^{N}\lambda^r=0$. These formulas follow first
for shifts in the right half-plane and then meromorphically. They
show that \eqref{ats:eq:polynomial} gives finite rational-color
Gamma/polygamma identities without appealing to numerical period
independence.

\subsection{The resolvent square and a two-shift extension}

\begin{corollary}[A sine-square evaluation]
\label{ats:cor:resolvent}
For $x\notin\mathbb Z\setminus\{0\}$,
\begin{equation}\label{ats:eq:resolvent}
 \boxed{\displaystyle
 \sum_{n\ge0}(2n+1)
   \left(\sum_{k>n}\frac{(-1)^k}{k^2-x^2}\right)^2
 =\frac{\pi^2}{4\sin^2(\pi x)}-\frac1{4x^2}.}
\end{equation}
The apparent singularity at zero is removable, with value
$\zeta(2)/2$.
\end{corollary}

\begin{proof}
Apply Theorem~\ref{ats:thm:polynomial} to $P(k)=k^2-x^2$ and
$\lambda=-1$. The cotangent/cosecant partial fractions give
\[
 E_0=\frac1{2x^2}-\frac{\pi}{2x\sin(\pi x)}.
\]
The endpoint is $x^2E_0^2-E_0$, which is the displayed expression.
The proof also establishes meromorphic continuation and local normal
convergence away from the stated coordinate poles.
\end{proof}

For example, at $x=1/2,1/3,1/4$ the values are, respectively,
\[
 \frac{\pi^2}{4}-1,\qquad
 \frac{\pi^2}{3}-\frac94,\qquad
 \frac{\pi^2}{2}-4.
\]
For real $y\ne0$, substituting $x=iy$ instead gives
\[
 \sum_{n\ge0}(2n+1)
   \left(\sum_{k>n}\frac{(-1)^k}{k^2+y^2}\right)^2
 =\frac1{4y^2}-\frac{\pi^2}{4\sinh^2(\pi y)}.
\]

If $T_s^-(n)=\sum_{k>n}(-1)^k k^{-s}$, coefficient extraction in
\eqref{ats:eq:resolvent} gives, for every integer $m\ge2$,
\begin{equation}\label{ats:eq:convolution}
 \sum_{a=1}^{m-1}\sum_{n\ge0}(2n+1)
            T_{2a}^-(n)T_{2m-2a}^-(n)
 =\frac{2m-3}{2}\zeta(2m-2).
\end{equation}
To check the coefficient, differentiate
$\pi\cot(\pi x)=x^{-1}-2\sum_{r\ge1}\zeta(2r)x^{2r-1}$.
The resolvent series is normally convergent for $|x|<1$, including
after any fixed coefficient extraction.

Write
\[
 \mathcal B(z)=\frac12\left\{\psi\left(\frac{z+1}{2}\right)
                              -\psi\left(\frac z2\right)\right\},
 \quad
 M_n(a,b)=\frac{\mathcal B(n+a+1)-\mathcal B(n+b+1)}{b-a}.
\]
The quotient has its derivative meaning at $a=b$.
Partial fractions and the alternating digamma series identify it with
$\sum_{j\ge0}(-1)^j/((j+n+a+1)(j+n+b+1))$.
Taking $P(k)=(k+a)(k+b)$ gives the all-shift identity
\begin{equation}\label{ats:eq:two-shift}
 \sum_{n\ge0}(2n+a+b+1)M_n(a,b)^2
 =M_0(a,b)-abM_0(a,b)^2.
\end{equation}
It is meromorphic in $a,b$, with the positive-lattice denominator
poles excluded. In particular, $a=0,b=1$ yields
\[
 \sum_{n\ge0}(2n+2)
  \left(\sum_{k>n}\frac{(-1)^k}{k(k+1)}\right)^2
 =2\log2-1.
\]

\subsection{Continuous order: an ordinary Hurwitz identity}

Define the alternating Hurwitz family
\begin{equation}\label{ats:eq:beta-s}
 \mathcal B_s(z)=\sum_{j\ge0}\frac{(-1)^j}{(j+z)^s}
 =2^{-s}\left\{\zeta\left(s,\frac z2\right)
              -\zeta\left(s,\frac{z+1}{2}\right)\right\}.
\end{equation}
The Hurwitz difference is entire in $s$, since its residues at one
cancel. On $\Re s>0,\Re z>0$ the ordinary alternating series and
\[
 \mathcal B_s(z)=\frac1{\Gamma(s)}
       \int_0^{\infty}\frac{u^{s-1}e^{-zu}}{1+e^{-u}}\,du
\]
agree. At $s=1$, this is the classical Nielsen beta function \cite{NielsenBeta}
$\mathcal B$; it is not Euler's two-variable beta function.

\begin{theorem}[Continuous-order square law]
\label{ats:thm:continuous}
For $\Re s>0$, $\Re a>0$, using the principal powers,
\begin{align}\label{ats:eq:continuous}
 &\sum_{n\ge0}\{(n+a+1)^s-(n+a)^s\}
                      \mathcal B_s(n+a+1)^2\notag\\
 &\hspace{18mm}=\mathcal B_s(a+1)-a^s\mathcal B_s(a+1)^2.
\end{align}
The series is absolutely and locally normally convergent in $(s,a)$.
The formula also holds at real $a=0$, with $0^s=0$ in this half-plane.
In that case the right side is the Dirichlet eta function
$\eta(s)=(1-2^{1-s})\zeta(s)$, with its removable value at one.
\end{theorem}

\begin{proof}
The recurrence
$\mathcal B_s(z)+\mathcal B_s(z+1)=z^{-s}$ implies, for
$F_s(z)=z^s\mathcal B_s(z)\mathcal B_s(z+1)$,
\[
 F_s(z-1)-F_s(z)=\{z^s-(z-1)^s\}\mathcal B_s(z)^2.
\]
Summing this finite identity gives the stronger statement
\begin{align}\label{ats:eq:continuous-finite}
 &\sum_{n=0}^{N}\{(n+a+1)^s-(n+a)^s\}
                          \mathcal B_s(n+a+1)^2\notag\\
 &=\mathcal B_s(a+1)-a^s\mathcal B_s(a+1)^2
 -(N+a+1)^s\mathcal B_s(N+a+1)\mathcal B_s(N+a+2).
\end{align}
The lower endpoint form follows from the same recurrence and does
not require evaluating $\mathcal B_s(0)$ at $a=0$.

The Laplace representation, or paired summation by parts, gives
$\mathcal B_s(n+a+1)=O(n^{-\Re s})$ locally uniformly. The weight
difference is $O(n^{\Re s-1})$, since it equals
$s\int_0^1(n+a+v)^{s-1}\,dv$ for large $n$.
The summands are therefore $O(n^{-\Re s-1})$, while the remaining
endpoint in \eqref{ats:eq:continuous-finite} is $O(N^{-\Re s})$.
This proves the asserted convergence and permits every fixed
spectral and shift derivative by Cauchy's formula.
\end{proof}

The integer specializations give
\[
 \sum_{n\ge0}\{(n+1)^d-n^d\}\mathcal B_d(n+1)^2=\eta(d)
 \qquad(d=1,2,\ldots).
\]
More generally, for $d\ge2$, the integrands and summands in
\eqref{ats:eq:continuous} are finite polygamma expressions because
\[
 \mathcal B_d(z)=
 \frac{(-1)^d}{2^d(d-1)!}
 \left\{\psi^{(d-1)}\left(\frac z2\right)
       -\psi^{(d-1)}\left(\frac{z+1}{2}\right)\right\}.
\]

\subsection{All Stieltjes orders with no finite-part prescription}

Use the Laurent convention
$\zeta(1+u,z)=u^{-1}+\sum_{j\ge0}(-1)^j\gamma_j(z)u^j/j!$.
Put $L=\log2$ and define
\begin{align*}
 C_r(n)&=\sum_{j=0}^{r}\binom rj L^{r-j}
 \left\{\gamma_j\left(\frac{n+1}{2}\right)
       -\gamma_j\left(\frac{n+2}{2}\right)\right\},\\
 W_k(n)&=(n+1)\log^k(n+1)-n\log^k n.
\end{align*}
At $n=0$, take $n\log^k n=0$ and $\log^0(n+1)=1$;
thus $W_0(n)=1$ for every $n$. Equation~\eqref{ats:eq:beta-s}
immediately gives
\[
 \left.\partial_s^r\mathcal B_s(n+1)\right|_{s=1}
       =\frac{(-1)^r}{2}C_r(n).
\]

\begin{corollary}[An all-index Stieltjes square identity]
\label{ats:cor:stieltjes}
For each integer $r\ge0$,
\begin{align}\label{ats:eq:stieltjes}
 &\frac14\sum_{n\ge0}
  \sum_{k+\ell+m=r}\frac{r!}{k!\ell!m!}
                 (-1)^k W_k(n)C_\ell(n)C_m(n)\notag\\
 &=\frac{L^{r+1}}{r+1}
          -\sum_{j=1}^{r}\binom rj L^j\gamma_{r-j}.
\end{align}
Every fixed summand family is absolutely convergent. The empty
right-hand sum at $r=0$ is zero.
\end{corollary}

\begin{proof}
Differentiate \eqref{ats:eq:continuous} $r$ times in $s$ at $a=0,s=1$.
Normal convergence has been proved on a larger compact neighborhood
of that point in the spectral half-plane. Leibniz's rule gives the
left side, after multiplication by $(-1)^r$.
Expanding $(1-e^{-Lu})\zeta(1+u)$ gives
\[
 \eta^{(r)}(1)=(-1)^r\left\{
  \frac{L^{r+1}}{r+1}
   -\sum_{j=1}^{r}\binom rj L^j\gamma_{r-j}\right\}.
\]
This is the required right side. More explicitly,
$C_j(n)=O(n^{-1}\log^j(n+1))$ and
$W_k(n)=O(\log^k(n+1))$, so all the displayed scalar series
converge absolutely. These estimates also follow by differentiating
the paired alternating series, keeping its cancellation.
\end{proof}

For instance $\sum_{n\ge0}C_0(n)^2=4\log2$, and the first
logarithmic identity is
\[
 \frac14\sum_{n\ge0}
  \{W_1(n)C_0(n)^2-2C_0(n)C_1(n)\}
 =\gamma\log2-\frac12\log^2 2.
\]
At second order the corresponding right side is
$\tfrac13\log^3 2-\gamma\log^2 2-2\gamma_1\log2$.
These are ordinary convergent identities on the half-lattice, not
distributional contact terms or finite parts.

\subsection{An anchored Gamma primitive and shift derivatives}

At $s=1$, \eqref{ats:eq:continuous} becomes
\begin{equation}\label{ats:eq:nielsen-square}
 \sum_{n\ge0}\mathcal B(n+a+1)^2
       =\mathcal B(a+1)-a\mathcal B(a+1)^2.
\end{equation}
For real $u\ge0$, its compensated primitive is
\begin{align}\label{ats:eq:gamma-primitive}
 &\int_0^u\left\{\sum_{n\ge0}\mathcal B(n+a+1)^2
                   +a\mathcal B(a+1)^2\right\}\,da\notag\\
 &=\log\frac{\sqrt\pi\,\Gamma(1+u/2)}{\Gamma((1+u)/2)}.
\end{align}
Every exchange follows from local normal convergence, and the
constant is fixed by $u=0$. The same identity holds holomorphically
on a simply connected pole-free neighborhood with the logarithm
anchored at zero. At $u=1$, the value is $\log(\pi/2)$.

Finally, because $\partial_a^j\mathcal B(n+a+1)=
(-1)^j j!\mathcal B_{j+1}(n+a+1)$, every shift derivative of
\eqref{ats:eq:nielsen-square} gives an explicit mixed-order identity.
Writing $z=a+1$ and
$U_r(z)=\sum_{j=0}^{r}\mathcal B_{j+1}(z)\mathcal B_{r-j+1}(z)$,
$U_{-1}=0$, one obtains for $r\ge0$
\begin{equation}\label{ats:eq:shift-convolution}
 \sum_{n\ge0}U_r(n+z)
 =\mathcal B_{r+1}(z)+U_{r-1}(z)-(z-1)U_r(z).
\end{equation}
At positive integer orders the right side is a finite product of
ordinary polygamma values at half-shifts.  This is a convolution
identity; it does not claim to isolate every individual mixed product.

\paragraph{Focused next questions.}
The diagonal finite certificate suggests a precise bilinear question:
for distinct polynomial denominators $P,Q$, classify rational-coefficient
certificates involving the two adjacent-tail products. A second target
is meromorphic continuation of the continuous-order square law past
$\Re s=0$, with the explicit upper endpoint in
\eqref{ats:eq:continuous-finite} retained at each resonance. A third is
to combine these matched-color identities with the independent
two-parameter conjugation generator, determining which mixed resolvent
coefficients are eliminated by their joint relation system. None of
these questions is a numerical-independence claim.

% END INCLUDED SOURCE: sections/02_tail_squares.tex

% BEGIN INCLUDED SOURCE: sections/03_cyclotomic_generator.tex
\section{A finite cyclotomic generator for products of rational tails}
\label{cyc:sec}

The incoming Shifts and Twists report \cite{HTST} asks for an analogue of its two-parameter
alternating identity at higher roots of unity.  There is a precise answer
for one conjugation sector, valid at every conductor.  The answer below
uses only elementary trigonometric functions and finitely many digamma
values.  Its coefficient form is consistent with the classical
depth-two parity theorem for colored multiple polylogarithms
\cite{PanzerParity}.  We give an independent bilateral-sum proof, including
the limiting conventions at the trivial color.  No reduction of the
opposite conjugation sector is asserted.

\subsection{The finite functions}

If $\omega$ is a root of unity, define, initially for $|t|<1$,
\begin{equation}\label{cyc:F}
 F_\omega(t)=\begin{cases}
 -\psi(1-t)-\gamma,&\omega=1,\\[2pt]
 -\dfrac1N\displaystyle\sum_{j=1}^N\omega^j
                  \psi\bigl((j-t)/N\bigr),
       &\omega^N=1,\quad\omega\ne1.
 \end{cases}
\end{equation}
The value is independent of the chosen multiple $N$ of the conductor.
For $\omega\ne1$ it is the ordinary convergent series
$\sum_{k\ge1}\omega^k/(k-t)$; residue-class grouping and the digamma
difference series prove \eqref{cyc:F}.  Its Taylor series is
\[
 F_\omega(t)=\sum_{j\ge1}\operatorname{Li}_j(\omega)t^{j-1}.
\]
For the trivial color the coefficient of $t^0$ in this last notation is
set to zero; all coefficients of positive degree are their ordinary
zeta values.  This convention does not assign a value to a divergent
harmonic series.

Write $\omega=e^{2\pi i\theta}$ with $0<\theta<1$ for a nontrivial
color and put
\begin{equation}\label{cyc:B}
 B_\omega(t)=\begin{cases}
 \pi t\cot(\pi t),&\omega=1,\\[2pt]
 \dfrac{\pi t e^{i\pi(2\theta-1)t}}{\sin(\pi t)},&\omega\ne1.
 \end{cases}
\end{equation}
These functions have value $1$ at zero.  They obey
\begin{equation}\label{cyc:BF}
 B_\omega(t)=1-tF_\omega(t)+tF_{\omega^{-1}}(-t).
\end{equation}
For example $B_{-1}(t)=\pi t/\sin(\pi t)$.
The coefficients $b_j(\omega)=[t^j]B_\omega(t)$ are explicit:
\[
 b_j(\omega)=\frac{(2\pi i)^j}{j!}B_j(\theta)\quad(\omega\ne1),
 \qquad b_0(1)=1,\quad b_1(1)=0,
 \quad b_j(1)=-2\zeta(j)\ (j\ge2\hbox{ even}),
\]
and $b_j(1)=0$ for odd $j\ge3$.  The $B_j(\theta)$ here are ordinary
Bernoulli polynomials, whereas $B_\omega(t)$ denotes their scaled
generating function.  Formula \eqref{cyc:BF} also proves that for $j\ge1$
\[
 b_j(\omega)=-\operatorname{Li}_j(\omega)
                    -(-1)^j\operatorname{Li}_j(\omega^{-1}),
\]
with the stated convention at $\omega=1$, $j=1$.

Define the following finite expression:
\begin{align}\label{cyc:C}
 C_{a,b}(X,Y)={}&
 \frac{B_{ab}(X)F_b(Y-X)-B_a(X)F_b(Y)}X\notag\\
 &-\frac{B_{ab}(Y)F_{a^{-1}}(Y-X)-F_{a^{-1}}(-X)}Y.
\end{align}
Each numerator has the indicated factor, so $X=0$ and $Y=0$ are
removable for this expression near the origin.  We use the finite
four-point operation
\begin{equation}\label{cyc:Q}
 Q_{a,b}(x,y)=\frac{
 C_{a,b}(x,y)+C_{a,b}(x,-y)
 -C_{a,b}(-x,y)-C_{a,b}(-x,-y)}{4x}.
\end{equation}
Finally set
\begin{equation}\label{cyc:J}
 J_\omega(t)=\frac{B_\omega(t)-B_\omega(-t)}{2t}
 =\begin{cases}
 0,&\omega=1,\\[2pt]
 \displaystyle i\pi\frac{\sin(\pi(2\theta-1)t)}{\sin(\pi t)},
       &\omega\ne1.
 \end{cases}
\end{equation}
The function $J_\omega$ is even.  Thus
$[J_\omega(x)-J_\omega(y)]/(x^2-y^2)$ has its usual derivative
interpretation when $y=\pm x$.

\subsection{All colors and an arbitrary outer character}

For roots of unity $a,b,c$ write
\begin{align}
 E_n^a(x)&=\sum_{k>n}\frac{a^k}{k^2-x^2},\notag\\
 R_{a,b;c}(x,y)&=\sum_{n\ge0}c^n\{c(n+1)-n\}
                                      E_n^a(x)E_n^b(y).
 \label{cyc:R}
\end{align}
The normalization of the weight makes its partial sum exactly
$c^m m$.  In particular $R_{a,b;-1}=-\mathcal K_{a,b}$ for the
alternating weight $\mathcal K_{a,b}=\sum(-1)^n(2n+1)E_n^aE_n^b$.

\begin{theorem}[Cyclotomic rational-tail generator]\label{cyc:master}
Let $a,b,c$ be arbitrary roots of unity and put $d=abc$.  Initially
assume $|x|+|y|<1$ and $xy\ne0$.  Then
\begin{align}\label{cyc:mastereq}
 &R_{a,b;c}(x,y)+R_{a^{-1},b^{-1};c^{-1}}(x,y)\notag\\
 &\qquad=Q_{a,bc}(x,y)+Q_{b,ac}(y,x)
                 +\frac{J_d(x)-J_d(y)}{x^2-y^2}.
\end{align}
Both sides extend meromorphically to $\mathbb C^2$.  All poles away
from $x\in\mathbb Z\setminus\{0\}$ or
$y\in\mathbb Z\setminus\{0\}$ are removable in the combined
right side.  The series is absolutely locally normally convergent
unless $a=b=1$ and $c\ne1$.  In that exceptional case it converges
ordinarily and locally uniformly; subtracting the summand at
$(x,y)=(0,0)$ makes the remaining series absolutely locally normally
convergent.  When $x,y$ are real, the left side of
\eqref{cyc:mastereq} is $2\Re R_{a,b;c}(x,y)$.
\end{theorem}

In particular this is a conductor-four generator, with quarter-shift
digamma expressions furnished by \eqref{cyc:F}.  It also works for
odd conductors, because the outer color $c$ is arbitrary.  The appearance
of a conjugate pair is essential: the theorem proves a depth reduction
for the real sector at odd total weight, and does not prove such a
reduction for each complex summand separately.

For nonzero rational $x,y$ with $x,y,x+y,x-y\notin\mathbb Z$,
Gauss's rational digamma formula \cite[Eq.~(5.4.19)]{DLMFGamma}
makes the answer a finite explicit
combination of $\pi$, $\pi^2$, and $\pi\log\alpha$ for positive
algebraic $\alpha$, with algebraic coefficients.  The same assertion
holds on $x=y\ne0$ when $2x\notin\mathbb Z$.  To see the stated
coefficient class, $F_\omega$ at a rational argument is an algebraic
linear combination of $1$, $\pi$, and such logarithms: its Euler
constant cancels in \eqref{cyc:F}.  Every $B_\omega$ value at a
nonzero rational argument is $\pi$ times an algebraic number.
The apparently unmultiplied term $F_{a^{-1}}(-X)/Y$ in
\eqref{cyc:C} is odd in $Y$ and disappears in \eqref{cyc:Q}.
The remaining $J$ quotient is elementary, also at its diagonal
limit.  This statement concerns the displayed nonsingular rational
parameters; the axes can introduce higher polylogarithmic constants,
as the $\zeta(3)$ limits below demonstrate.

\subsection{A bilateral proof with finite cutoff control}

For nontrivial unit colors $a,b$, introduce the ordinary generating
series
\[
 Z_{a,b}(X,Y)=\sum_{k>m\ge1}\frac{a^k b^m}{(k-X)(m-Y)}.
\]
Summation by parts in $k$ proves locally uniform convergence: the inner
partial sums grow at most logarithmically and their successive
increments are $O(1/k)$.  We claim
\begin{equation}\label{cyc:bilateral-generator}
 Z_{a,b}(X,Y)-Z_{a^{-1},b^{-1}}(-X,-Y)
 =C_{a,b}(X,Y)-\frac{F_{ab}(X)-F_{ab}(Y)}{X-Y}.
\end{equation}
The last quotient is removable at $X=Y$.

To prove it, sum first over the full integer lattice in
\[
 S=\sum_{m\ge1}b^m\sum_{k\in\mathbb Z}
       \frac{(ab)^k}{(k-X)(k+m-Y)}.
\]
For each fixed $m$ the inner series is absolutely convergent.  The
bilateral partial-fraction formula is
\[
 \operatorname{PV}\sum_{k\in\mathbb Z}\frac{\omega^k}{k-t}
                  =-\frac{B_\omega(t)}t.
\]
It follows either from \eqref{cyc:BF} or from the elementary
Fourier series of the shifted resolvent.  The symmetric principal
value at $\omega=1$ is exactly the cotangent in \eqref{cyc:B}.
Partial fractions therefore give
\begin{equation}\label{cyc:S-eval}
 S=-\frac{B_{ab}(X)}X F_b(Y-X)
                  +\frac{B_{ab}(Y)}YF_{a^{-1}}(Y-X).
\end{equation}
Now put $l=k+m$.  The positive, negative, mixed, and two zero-index
parts of the region $l>k$ give respectively
\begin{align*}
 S={}&Z_{b,a}(Y,X)+Z_{a^{-1},b^{-1}}(-X,-Y)
       -F_{a^{-1}}(-X)F_b(Y)\\
    &-F_b(Y)/X+F_{a^{-1}}(-X)/Y.
\end{align*}
This rearrangement has a direct cutoff justification.  Keep
$1\le m\le M$ before separating regions.  In each same-sign region,
the omitted strip is bounded after summation by parts by a constant
times
\[
 \sum_{k\ge1}\frac1{k(k+M)}=O(\log M/M).
\]
In the mixed region the difference between the triangular cutoff
$u+v\le M$ and the product of the two ordinary sums is
$O(\log M/M)$ as well: for $u<M$ use the nontrivial-color tail
$O((M-u)^{-1})$, and for $u\ge M$ use the nontrivial-color tail in
$u$.  The two zero-index tails are $O(1/M)$.  These estimates are
uniform on compact parameter sets avoiding poles.  Thus all the
displayed limiting rearrangements are valid.

The elementary product decomposition is
\[
 Z_{b,a}(Y,X)=F_a(X)F_b(Y)-Z_{a,b}(X,Y)
                    -\frac{F_{ab}(X)-F_{ab}(Y)}{X-Y}.
\]
Use this identity, \eqref{cyc:BF}, and \eqref{cyc:S-eval} to obtain
\eqref{cyc:bilateral-generator}.

Here are explicit conventions for the boundary colors.  For $r\ge2$,
$s\ge1$, let $w=r+s$ and write $L_j(\omega)=\operatorname{Li}_j(\omega)$.
Coefficient extraction in \eqref{cyc:bilateral-generator} gives
\begin{align}\label{cyc:finite-parity}
 &Z(r,s;a,b)-(-1)^w Z(r,s;a^{-1},b^{-1})\notag\\
 ={}&(-1)^s\sum_{j=s}^{w}(-1)^j
          \binom{j-1}{s-1}L_j(b)b_{w-j}(ab)-L_w(ab)\notag\\
 &+(-1)^r\sum_{j=r}^{w}\binom{j-1}{r-1}
               L_j(a^{-1})b_{w-j}(ab)-L_s(b)b_r(a),
\end{align}
where
$Z(r,s;a,b)=\sum_{k>m\ge1}a^k b^m/(k^r m^s)$.
Every double series here is absolutely convergent.  If $b=1$, $s=1$,
the first summand with $j=1$ cancels the last term before any value
is assigned to $L_1(1)$.  If $ab=1$, the coefficient of $b_1(ab)$
vanishes, since $b=a^{-1}$ and
$\binom{w-2}{s-1}=\binom{w-2}{r-1}$ with opposite signs.
There is no other $L_1(1)$ term because $r\ge2$.  Consequently the
finite formula, and its normally convergent even-$r$, odd-$s$
generator, extend to every root-of-unity color with exactly
\eqref{cyc:F}--\eqref{cyc:B}.  One may also obtain this extension
by taking colors to $1$ on the unit circle; the only discarded
terms have $r=1$, which the projection in \eqref{cyc:Q} annihilates.

We now prove the tail theorem.  A finite outer cutoff, followed by
the telescoping identity
\[
 \sum_{n=0}^{m-1}c^n\{c(n+1)-n\}=c^m m,
\]
gives
\begin{equation}\label{cyc:ordered-tail}
 R_{a,b;c}(x,y)=D_{a,bc}(x,y)+D_{b,ac}(y,x)+H_{abc}(x,y),
\end{equation}
where
\begin{align*}
 D_{u,v}(x,y)&=\sum_{k>m\ge1}
   \frac{u^k v^m m}{(k^2-x^2)(m^2-y^2)},\\
 H_d(x,y)&=\frac{O_d(x)-O_d(y)}{x^2-y^2},\qquad
 O_d(t)=\frac{F_d(t)+F_d(-t)}2.
\end{align*}
The double sums are absolutely convergent, dominated by a constant
times $\sum_{k>m}k^{-2}m^{-1}$.  The omitted cutoff boundary is
$O(M^{-1})$, even when both tails are untwisted.  Applying the
four-point operation to \eqref{cyc:bilateral-generator} gives
\[
 D_{u,v}+D_{u^{-1},v^{-1}}=Q_{u,v}-H_{uv}.
\]
Add \eqref{cyc:ordered-tail} to its inverse-color version.  The two
copies of $-H_d$, together with the two diagonal terms, leave
$H_{d^{-1}}-H_d$.  Formula \eqref{cyc:BF} says
$O_{d^{-1}}-O_d=J_d$, proving \eqref{cyc:mastereq}.

For completeness, locally uniformly away from the coordinate poles,
\[
 E_n^1(x)=O(n^{-1}),\qquad E_n^a(x)=O(n^{-2})\quad(a\ne1),
 \qquad E_n^a(x)-E_n^a(0)=O(n^{-3}).
\]
These imply every asserted absolute convergence statement.  If
$a=b=1$ and $c\ne1$, the scalar summand at zero is
$c^n\{(c-1)/n+O(n^{-2})\}$ for $n\ge1$, which is ordinarily
convergent; the parameter difference is $O(n^{-3})$.  The case
$c=1$ has weight $1$ and is absolutely convergent even for $a=b=1$.
Near a coordinate pole only finitely many tail terms are singular.
Normal convergence and the meromorphic identity theorem therefore
give the claimed continuation and all oblique removability.

\subsection{Higher polynomial weights and their precise parity sector}

For $p,q\ge2$, $1\le h<\min(p,q)$, put
\[
 T_p^a(n)=\sum_{k>n}\frac{a^k}{k^p},\qquad
 M_{p,q,h}(a,b;c)=\sum_{n\ge0}c^n
       \{c(n+1)^h-n^h\}T_p^a(n)T_q^b(n).
\]
Exactly the same cutoff proof gives the ordinary identity
\begin{equation}\label{cyc:allweights}
 M_{p,q,h}(a,b;c)=Z(p,q-h;a,bc)+Z(q,p-h;b,ac)
                         +L_{p+q-h}(abc).
\end{equation}
The only possibly conditional case is $p=q=2$, $h=1$,
$a=b=1$, $c\ne1$; when $c=1$ the weight has one lower degree.
Let $w=p+q-h$.  Formula \eqref{cyc:finite-parity} now explicitly
reduces
\[
 M_{p,q,h}(a,b;c)-(-1)^w
 M_{p,q,h}(a^{-1},b^{-1};c^{-1})
\]
to finitely many single colored polylogarithm values and Bernoulli
coefficients.  For real parameters this is the real sector if $w$
is odd and the imaginary sector if $w$ is even.  Thus the parity
qualification is fully explicit at every conductor and every
admissible weight.

\subsection{Two conductor-four diagonal identities}

Write
\[
 O_-(x)=-\log2+O_+(x/2)-O_+(x),\qquad
 O_+(x)=-\frac{\psi(1+x)+\psi(1-x)+2\gamma}{2}.
\]
For real $x$ away from nonzero integers, define
\[
 D(x)=\Re\sum_{n\ge0}(-1)^n(2n+1)E_n^i(x)^2,\qquad
 V(x)=\sum_{n\ge0}(-1)^n(2n+1)|E_n^i(x)|^2.
\]
Both are ordinary absolutely convergent sums.  The all-color theorem
gives the compact evaluations
\begin{align}
 D(x)&=\frac{\pi}{2x\sin(\pi x)}
   \{\cos(\pi x/2)O_-(x/2)-\cos(\pi x)O_-(x)\},
 \label{cyc:Ddiag}\\
 V(x)&=\frac{\pi}{2x\sin(\pi x)}
   \{\cos(\pi x/2)O_-(x/2)-O_-(x)\}.
 \label{cyc:Vdiag}
\end{align}
One verification of the simplification is to use
\[
 O_i(t)+O_{-i}(t)=O_-(t/2),\quad
 \frac{B_i(t)+B_i(-t)}2
       =\frac{\pi t\cos(\pi t/2)}{\sin(\pi t)},
 \quad B_1(t)=\pi t\cot(\pi t)
\]
in \eqref{cyc:C}--\eqref{cyc:Q}.  At $a=b=i$ the product color in
the ordered sum is $1$, while at $a=i,b=-i$ it is $-1$;
in either case $J$ vanishes.  The terms at $F_{\pm i}(0)$ cancel
in the conjugate sum.  This gives \eqref{cyc:Ddiag} and
\eqref{cyc:Vdiag} directly, without a limiting singular quotient.

In particular, subtracting these two formulas proves the especially
short sine-tail square
\begin{equation}\label{cyc:sinesquare}
 \boxed{\displaystyle
 \sum_{n\ge0}(-1)^n(2n+1)
  \left(\sum_{k>n}\frac{\sin(\pi k/2)}{k^2-x^2}\right)^2
       =-\frac{\pi}{4x}\tan(\pi x/2)O_-(x).}
\end{equation}
The algebraic-square version extends meromorphically to complex $x$;
no complex conjugation is inserted there.  Its removable value at
zero is $\pi^2\log2/8$.  Two rational values are
\begin{align}
 \sum_{n\ge0}(-1)^n(2n+1)
  \left(\sum_{k>n}\frac{\sin(\pi k/2)}{k^2-1/4}\right)^2
       &=\frac\pi2,\label{cyc:sinehalf}\\
 \sum_{n\ge0}(-1)^n(2n+1)
  \left(\sum_{k>n}\frac{\sin(\pi k/2)}{k^2-1/16}\right)^2
       &=\pi(\sqrt2-1)
                 \{2-\sqrt2\log(1+\sqrt2)\}.
 \label{cyc:sinequarter}
\end{align}
Here $O_-(1/2)=-1$ and
$O_-(1/4)=\sqrt2\log(1+\sqrt2)-2$, by the rational digamma
formula.  The same substitution gives
\[
 D(1/2)=\pi\{\log(1+\sqrt2)-\sqrt2\},\qquad
 V(1/2)=\pi\{1+\log(1+\sqrt2)-\sqrt2\}.
\]
The coefficient limits are
\[
 D(0)=\frac9{32}\zeta(3)-\frac{3\pi^2\log2}{16},\qquad
 V(0)=\frac9{32}\zeta(3)+\frac{\pi^2\log2}{16}.
\]
For the cosine tails, the corresponding square has value
$9\zeta(3)/32-\pi^2\log2/16$ at zero.  These are evaluations
of an ordinary alternating outer sum, so a negative value of the
cosine square causes no positivity issue.

\subsection{An anchored quarter-color Gamma primitive}

Let $S(x)$ denote the holomorphic function in
\eqref{cyc:sinesquare} on $|x|<1$.  Then
\begin{equation}\label{cyc:gamma-primitive}
 \boxed{\displaystyle
 \int_0^u x\cot(\pi x/2)S(x)\,dx
  =\frac\pi8\log\!\left(
  \frac{\Gamma(1+u/2)\Gamma((1-u)/2)}
       {\Gamma(1-u/2)\Gamma((1+u)/2)}\right),\qquad |u|<1.}
\end{equation}
The integrand has a removable singularity at zero.  The logarithm
is its holomorphic branch equal to zero at $u=0$, and the path can
be any path in the unit disk.  Indeed \eqref{cyc:sinesquare} makes
the integrand $-\pi O_-(x)/4$.  Integrating the displayed digamma
formula for $O_-$ and then applying Gamma duplication yields
\eqref{cyc:gamma-primitive}; the two sides vanish at zero.
Absolute local normal convergence permits integration term by term
in the defining sum for $S$.  A concrete endpoint is
\[
 \int_0^{1/2}x\cot(\pi x/2)S(x)\,dx
       =\frac\pi4\log\frac{\Gamma(1/4)}{2\Gamma(3/4)}.
\]
This primitive is a direct conductor-four continuation of the
earlier alternating diagonal Gamma primitive, with its additive
constant fixed explicitly.

% END INCLUDED SOURCE: sections/03_cyclotomic_generator.tex

% BEGIN INCLUDED SOURCE: sections/04_spectral_defects.tex
\section{Every spectral derivative of the conjugate-factor defect}
\label{sd:section}

The quadratic-spectrum report \cite{QSpec} proves a first-derivative
factorization law and asks, in its second further-research question, for
the higher-derivative correction.  The answer below consists of an
ordinary convergent logarithmic sum and finitely many explicitly stated
Hurwitz or Stieltjes terms.  It works for every polynomial weight and
every positive mass.  The finite term arising from one extra power of
the spectral variable is essential.

\subsection{Conventions and the classical first derivative}

Let $a,b>0$, put $x_n=n+a$ and $Q_b(x)=x^2+b^2$, and let
$W(x)=\sum_{d=0}^{D}w_dx^d\in\mathbb C[x]$.  Initially in their
half-planes of absolute convergence, define
\begin{align}
 F_{W,\pm}(s)&=\sum_{n\ge0}W(x_n)(x_n\pm ib)^{-s},
 &Z_W(s)&=\sum_{n\ge0}W(x_n)Q_b(x_n)^{-s}.
 \label{sd:def:spectra}
\end{align}
The powers use principal logarithms.  Since $x_n>0$,
$\log(x_n+ib)+\log(x_n-ib)=\log Q_b(x_n)$ without a branch correction.
All evaluations at zero refer to meromorphic continuation.  We use the
normalization requested in \cite{QSpec}:
\begin{align}
 \mathscr D_W(s)&=\frac{F_{W,+}(2s)+F_{W,-}(2s)}2-Z_W(s),
 \label{sd:def:difference}\\
 \Delta_r(W;a,b)&=\mathscr D_W^{(r)}(0)
 =2^{r-1}\{F_{W,+}^{(r)}(0)+F_{W,-}^{(r)}(0)\}-Z_W^{(r)}(0).
 \label{sd:def:defect}
\end{align}
Primes and parenthesized derivatives on these spectra always
differentiate their order, not their shift or mass.

For reference, the established local anomaly is
\begin{align}
 \Delta_1(W;a,b)
 &=\frac14[x^{-1}]W(x)
       \left\{\log(1+ib/x)-\log(1-ib/x)\right\}^{2}
       \label{sd:eq:first-residue}\\
 &=\sum_{k\ge1}w_{2k-1}\frac{(-1)^kb^{2k}}{k}
       \left(H_{2k}-\frac12H_k\right).
       \label{sd:eq:first-harmonic}
\end{align}
Only finitely many terms occur.  Its independence of $a$ and its
vanishing for every $b$ precisely when $W$ is even are inherited
results, not new claims of this article.  The general local-anomaly
framework is classical; see Dowker \cite{Dowker}.  We recover the formula
below to fix the normalization and expose the additional terms at
higher orders.

\subsection{Finite harmonic coefficients and global subtraction}

For $k\ge1$ define the polynomial
\begin{equation}\label{sd:eq:Ck}
 C_k(s;b)=(-1)^kb^{2k}
 \left\{\frac{(2s)_{2k}}{(2k)!}-\frac{(s)_k}{k!}\right\}
 =\sum_{m=2}^{2k}\alpha_{k,m}(b)s^m.
\end{equation}
Here $(s)_k=s(s+1)\cdots(s+k-1)$.  For a finite list $v$,
$e_j(v)$ denotes its elementary symmetric polynomial, with $e_0=1$
and $e_j=0$ beyond the length of the list.  Set
\begin{equation}\label{sd:eq:elementary}
 E_j(N)=e_j(1,1/2,\ldots,1/N),\qquad E_j(0)=\begin{cases}1&j=0,\\0&j>0.\end{cases}
\end{equation}
These are polynomials in the generalized harmonic numbers:
\[
 \sum_{j\ge0}E_j(N)t^j
 =\exp\left(\sum_{\ell\ge1}
       \frac{(-1)^{\ell-1}}{\ell}H_N^{(\ell)}t^\ell\right).
\]

\begin{lemma}[All local coefficients]\label{sd:lem:coefficients}
For $m\ge2$,
\begin{equation}\label{sd:eq:alpha}
 \alpha_{k,m}(b)=\frac{(-1)^kb^{2k}}{k}
       \{2^{m-1}E_{m-1}(2k-1)-E_{m-1}(k-1)\}.
\end{equation}
Set $\alpha_{k,m}=0$ for $m<2$ or $m>2k$.  In particular,
\[
 C_1(s;b)=-2b^2s^2,\qquad
 \alpha_{k,2}=\frac{(-1)^kb^{2k}}k
                    (2H_{2k-1}-H_{k-1}).
\]
\end{lemma}
\begin{proof}
Use
\[
 \frac{(s)_k}{k!}=\frac{s}{k}
                    \prod_{j=1}^{k-1}(1+s/j),\qquad
 \frac{(2s)_{2k}}{(2k)!}=\frac{s}{k}
                    \prod_{j=1}^{2k-1}(1+2s/j).
\]
Their linear terms agree.  Comparing every subsequent coefficient
proves the assertion, including the stated finite support.
\end{proof}

Put
\begin{equation}\label{sd:eq:h}
 h(s,x)=\frac{(x+ib)^{-2s}+(x-ib)^{-2s}}2-Q_b(x)^{-s},
 \qquad K=\left\lfloor\frac{D+1}{2}\right\rfloor.
\end{equation}
For $D=0$ the integer $K$ is zero and all sums from $1$ to $K$ are empty.

\begin{lemma}[A globally valid analytic subtraction]
\label{sd:lem:subtraction}
In a complex neighborhood of $s=0$,
\begin{align}
 \mathscr D_W(s)
 &=\sum_{n\ge0}W(x_n)
 \left\{h(s,x_n)-x_n^{-2s}
                  \sum_{k=1}^{K}C_k(s;b)x_n^{-2k}\right\}
 \notag\\
 &\quad+\sum_{d=0}^{D}\sum_{k=1}^{K}
         w_dC_k(s;b)\zeta(2s+2k-d,a).
 \label{sd:eq:global-subtraction}
\end{align}
The infinite sum converges locally normally.  All the apparent
singularities at zero in the finite sum are removable.  The original
three spectra in \eqref{sd:def:spectra} are holomorphic at zero.
\end{lemma}
\begin{proof}
For $x>b$, the ordinary binomial expansions give
\[
 h(s,x)=x^{-2s}\sum_{k\ge1}C_k(s;b)x^{-2k}.
\]
On compact mass sets, its remainder after $K$ terms is uniformly
$O(x^{-2\Re s-2K-2})$ on a small closed $s$-disk.  After multiplication
by $W$, this is $O(x^{-2+2\varepsilon})$ when $D$ is even and
$O(x^{-3+2\varepsilon})$ when $D$ is odd, for $|s|\le\varepsilon<1/4$.
This proves normal convergence.  The finitely many $x\le b$ use the
exact expression $h$, so no assumption $b<a$ is needed.

In a common right half-plane every term being subtracted is separately
summable; adding it back as a Hurwitz zeta proves
\eqref{sd:eq:global-subtraction} there and hence by continuation.
The only possible singular terms near zero have $2k-d=1$.
The simple pole of $\zeta(1+2s,a)$ is multiplied by $C_k=O(s^2)$,
so they are removable.  Finally, expansion of the polynomial weight
about $x\pm ib$ expresses each $F_{W,\pm}$ as a finite sum of
$\zeta(s-j,a\pm ib)$ with $j\ge0$.  They are holomorphic at zero,
and the same follows for $Z_W$ from \eqref{sd:def:difference}.
\end{proof}

\subsection{The complete higher-derivative formula}

Define the elementary logarithmic polynomial
\begin{align}
 H_r(x;b)
 &=(-1)^r\left[2^{r-1}
       \{\log^r(x+ib)+\log^r(x-ib)\}-\log^r Q_b(x)\right]
       \label{sd:eq:H-log}\\
 &=\sum_{j=1}^{\lfloor r/2\rfloor}
       (-1)^{r+j}4^j\binom r{2j}
       \log^{r-2j}Q_b(x)\arctan^{2j}(b/x),
       \label{sd:eq:H-real}
\end{align}
so $H_1=0$, $H_2=-4\arctan^2(b/x)$, and
$H_3=12\log Q_b(x)\arctan^2(b/x)$.  Also put
\begin{equation}\label{sd:eq:Pr}
 P_{k,r}(L)=r!\sum_{m=2}^{\min(r,2k)}
          \alpha_{k,m}(b)\frac{(-2L)^{r-m}}{(r-m)!}.
\end{equation}

\begin{theorem}[Every order, every polynomial weight, every positive mass]
\label{sd:thm:all-orders}
For every integer $r\ge1$,
\begin{align}
 \Delta_r(W;a,b)
 &=\sum_{n\ge0}W(x_n)
       \left\{H_r(x_n;b)-\sum_{k=1}^{K}
                    \frac{P_{k,r}(\log x_n)}{x_n^{2k}}\right\}
       \notag\\
 &\quad+\sum_{d=0}^{D}\sum_{k=1}^{K}w_d\mathcal J_{k,r}(2k-d;a,b),
       \label{sd:eq:all-orders}
\end{align}
where, if $\nu\ne1$,
\begin{equation}\label{sd:eq:nonresonant}
 \mathcal J_{k,r}(\nu;a,b)=r!\sum_{m=2}^{\min(r,2k)}
   \alpha_{k,m}(b)\frac{2^{r-m}}{(r-m)!}
                       \zeta^{(r-m)}(\nu,a),
\end{equation}
and at the resonant value $\nu=1$,
\begin{equation}\label{sd:eq:resonant}
 \boxed{\displaystyle
 \mathcal J_{k,r}(1;a,b)=r!\left\{
   \frac{\alpha_{k,r+1}(b)}2+
   \sum_{m=2}^{\min(r,2k)}\alpha_{k,m}(b)
          \frac{(-2)^{r-m}}{(r-m)!}\gamma_{r-m}(a)\right\}.}
\end{equation}
The infinite sum in \eqref{sd:eq:all-orders} is ordinary and absolutely
convergent.  It and every fixed parameter derivative converge locally
normally on positive parameter sets.
\end{theorem}
\begin{proof}
Differentiate the normally convergent expression
\eqref{sd:eq:global-subtraction}.  The derivative of its exact pointwise
term is $H_r$; writing
$\log(x\pm ib)=\frac12\log Q_b(x)\pm i\arctan(b/x)$ proves
\eqref{sd:eq:H-real}.  Differentiation of $x^{-2s}C_k(s;b)$ gives
\eqref{sd:eq:Pr}.  If $\nu\ne1$, Taylor multiplication gives
\eqref{sd:eq:nonresonant}.

At the pole, use the stipulated Laurent convention:
\[
 \zeta(1+2s,a)=\frac1{2s}
       +\sum_{j\ge0}\frac{(-2)^j}{j!}\gamma_j(a)s^j.
\]
The coefficient of $s^r$ in $C_k(s;b)\zeta(1+2s,a)$ contains
$\alpha_{k,r+1}/2$ from the pole and the displayed finite convolution
from its regular part.  Multiplication by $r!$ proves
\eqref{sd:eq:resonant}.  The uniform estimates in the lemma allow
the differentiations; logarithmic factors from spectral derivatives
do not affect summability.  The same estimates on small complex
neighborhoods of positive $(a,b)$ justify parameter derivatives.
\end{proof}

At $r=1$ the ordinary sum and every nonresonant term vanish.  The
remaining term is $\alpha_{k,2}/2$, which recovers
\eqref{sd:eq:first-harmonic}.  At $r\ge2$, deleting the ordinary sum
is generally incorrect.  Deleting $\alpha_{k,r+1}/2$ is also
incorrect when this coefficient is present: differentiating a regular
part alone loses the contribution of the pole.

\subsection{A single monomial isolates every Stieltjes index}

The simplest odd weight makes the correction especially transparent.
For $r\ge2$ set
\[
 c_r=(-1)^{r-1}2^{r-1}r(r-1).
\]
Since $C_1=-2b^2s^2$, the general theorem gives
\begin{equation}\label{sd:eq:linear-weight}
 \boxed{\displaystyle
 \Delta_r(x;a,b)=
 \sum_{n\ge0}\left\{x_nH_r(x_n;b)
       -c_rb^2\frac{\log^{r-2}x_n}{x_n}\right\}
       +c_rb^2\gamma_{r-2}(a).}
\end{equation}
Its residual summands are
$O(x_n^{-3}(1+\log^{r-2}x_n))$.  Thus this is an ordinary convergent
identity for each generalized Stieltjes constant, coupled to precisely
specified spectral derivatives.  For example,
\begin{align}
 \Delta_2(x;a,b)
 &=-4\sum_{n\ge0}
       \left\{x_n\arctan^2(b/x_n)-\frac{b^2}{x_n}\right\}
       +4b^2\psi(a),\label{sd:eq:linear-second}\\
 \Delta_3(x;a,b)
 &=12\sum_{n\ge0}
       \left\{x_n\log Q_b(x_n)\arctan^2(b/x_n)
                      -\frac{2b^2\log x_n}{x_n}\right\}
       +24b^2\gamma_1(a).\label{sd:eq:linear-third}
\end{align}
At first order, $\Delta_1(x;a,b)=-b^2$ even though each ordinary
pointwise first-derivative difference is zero.  This is a useful exact
check on the chosen continuation, not a license to omit higher-order
ordinary terms.

\subsection{The first-derivative parity criterion fails completely at order two}

\begin{theorem}[Universal second-defect vanishing]
\label{sd:thm:vanishing}
Fix $a>0$.  For a polynomial $W\in\mathbb C[x]$,
\[
 \Delta_2(W;a,b)=0\ \text{for every }b>0
 \quad\Longleftrightarrow\quad W=0.
\]
Thus no nonzero even polynomial retains the universal vanishing
property of the first derivative.
\end{theorem}
\begin{proof}
For $0<b<a$, the full binomial expansion converges normally.  If
$k>(D+1)/2$, all the Hurwitz values $\zeta(2k-d,a)$ in its
$b^{2k}$ coefficient are ordinary convergent sums.  That coefficient
in $\Delta_2$ is
\begin{equation}\label{sd:eq:high-mass-coeff}
 \frac{2(-1)^k}{k}(2H_{2k-1}-H_{k-1})
             \sum_{n\ge0}\frac{W(n+a)}{(n+a)^{2k}}.
\end{equation}
The scalar is nonzero.  If the defect vanishes identically, all these
moments vanish.  If $W\ne0$, let $n_0$ be the least nonnegative integer
with $W(n_0+a)\ne0$.  Multiplying the vanishing moment by
$(n_0+a)^{2k}$ and letting $k\to\infty$ gives
$W(n_0+a)=0$, a contradiction.  Dominated convergence is valid after
fixing any sufficiently large $k_0$: polynomial growth of $W$ makes
$\sum_{n>n_0}|W(n+a)|((n_0+a)/(n+a))^{2k_0}$ finite.
The reverse implication is immediate.
\end{proof}

This is a statement about identity in the mass at a fixed shift.
It does not exclude accidental zeros at a prescribed mass for signed
polynomial weights.  No inequality or numerical independence assertion
is needed for the proof.


% END INCLUDED SOURCE: sections/04_spectral_defects.tex

% BEGIN INCLUDED SOURCE: sections/05_gamma_mass_primitives.tex
\section{The logarithmic correction as a Gamma integral and its primitives}
\label{gp:section}

For the constant weight no subtraction is necessary at any positive
spectral derivative order.  The second derivative produces a particularly
simple ordinary function,
\begin{equation}\label{gp:eq:S}
 S(a,b)=\sum_{n\ge0}\arctan^2\!\left(\frac{b}{n+a}\right),
 \qquad a,b>0.
\end{equation}
Theorem~\ref{sd:thm:all-orders} immediately gives
\begin{equation}\label{gp:eq:zeta-second}
 Z_1''(0;a,b)=2\zeta''(0,a+ib)+2\zeta''(0,a-ib)+4S(a,b).
\end{equation}
The prime on $\zeta$ denotes its order derivative.  Equation
\eqref{gp:eq:zeta-second} is an exact higher-derivative replacement
for a factorization rule, not a claim that the extra function is
algebraically independent of some larger special-function class.

\begin{theorem}[An ordinary logarithmic Gamma integral]
\label{gp:thm:gamma-integral}
For $a,b>0$,
\begin{equation}\label{gp:eq:gamma-integral}
 \boxed{\displaystyle
 S(a,b)=2\int_0^1
  \frac{\log|\Gamma(a+ibt)|-\log|\Gamma(a+ib)|}{1-t^2}\,dt.}
\end{equation}
The integrand has a removable value at $t=1$.  Both sides are real
analytic in the positive parameters.
\end{theorem}
\begin{proof}
The Weierstrass product for the Gamma function, or the integrated
digamma series, gives the ordinary normally convergent identity
\[
 \log|\Gamma(a+ibt)|-\log|\Gamma(a+ib)|
 =\frac12\sum_{n\ge0}
          \log\frac{(n+a)^2+b^2}{(n+a)^2+b^2t^2}.
\]
For $0\le t\le1$ its summands are nonnegative, so Tonelli's theorem
permits the required integral.  The elementary identity
\begin{equation}\label{gp:eq:elementary-kernel}
 \int_0^1\frac{\log((1+y^2)/(1+y^2t^2))}{1-t^2}\,dt
 =\arctan^2y\qquad(y\ge0)
\end{equation}
finishes the proof.  To verify it, differentiate in $y$; cancellation
of $1-t^2$ leaves
$2y(1+y^2)^{-1}\int_0^1(1+y^2t^2)^{-1}dt
=2\arctan y/(1+y^2)$, and both sides vanish at zero.
The endpoint limit of the full integration density in
\eqref{gp:eq:gamma-integral}, including the prefactor $2$, is
$b\,\Im\psi(a+ib)$.  The defining series and its differentiated
tails converge uniformly on compact positive parameter sets.
\end{proof}

Gamma reflection gives two elementary-kernel specializations:
\begin{align}
 S(1,b)&=\int_0^1
       \frac{\log\{t\sinh(\pi b)/\sinh(\pi bt)\}}{1-t^2}\,dt,
       \label{gp:eq:integer}\\
 S(1/2,b)&=\int_0^1
       \frac{\log\{\cosh(\pi b)/\cosh(\pi bt)\}}{1-t^2}\,dt.
       \label{gp:eq:half}
\end{align}
The endpoint at $t=0$ in the first line is understood by continuity.
We do not assert an unproved finite classical-polylogarithm evaluation
of these integrals at generic $b$.

The small-mass expansion is an exact harmonic--Hurwitz generating function:
\begin{equation}\label{gp:eq:smallmass}
 S(a,b)=\sum_{k\ge1}\frac{(-1)^{k+1}b^{2k}}{k}
                 \left(H_{2k}-\frac12H_k\right)\zeta(2k,a),
 \qquad |b|<a.
\end{equation}
It follows by expanding $\arctan^2y$, or from
\eqref{sd:eq:high-mass-coeff}.  For every positive $b$, remove a
finite head with $N+a>b$ and replace $\zeta(2k,a)$ in the tail by
$\zeta(2k,N+a)$.  This gives a convergent implementation without
extending the disk of \eqref{gp:eq:smallmass} improperly.

Two useful exact transport rules are
\begin{align}
 S(a,b)-S(a+1,b)&=\arctan^2(b/a),\label{gp:eq:shift}\\
 \sum_{j=0}^{q-1}S\left(\frac{a+j}{q},\frac bq\right)&=S(a,b)
 \qquad(q\in\mathbb N).\label{gp:eq:distribution}
\end{align}
They follow by taking out the first term and by partitioning the
positive shifted lattice into residue classes.  The Gamma-integral
formula gives the same distribution law through Gauss multiplication.

\subsection{Every anchored shift antiderivative}

Let $H_0=0$ and, for $m\ge0$, define
\begin{equation}\label{gp:eq:J}
 J_m(z)=\frac{\zeta'(-m,z)+H_m\zeta(-m,z)}{m!},\qquad\Re z>0.
\end{equation}
The combination is chosen to remove the extra term generated by
differentiating the Pochhammer symbol.

\begin{lemma}\label{gp:lem:J}
For $m\ge1$,
\[
 \frac{d^m}{dz^m}J_m(z)=\zeta'(0,z)
             =\log\Gamma(z)-\frac12\log(2\pi).
\]
\end{lemma}
\begin{proof}
Differentiate the classical identity
$\partial_z^m\zeta(s,z)=(-1)^m(s)_m\zeta(s+m,z)$ in $s$ and set
$s=-m$.  Here $(s)_m=(-1)^mm!$ and its logarithmic derivative is
$-H_m$.  Thus
$\partial_z^m\zeta'(-m,z)=m!\{\zeta'(0,z)-H_m\zeta(0,z)\}$,
while $\partial_z^m\zeta(-m,z)=m!\zeta(0,z)$.  The terms cancel as
claimed.  The final equality is the standard Hurwitz--Gamma formula
\cite{DLMFZeta}.
\end{proof}

Define
\begin{equation}\label{gp:eq:I}
 I_m(a,b)=2\Re\int_0^1
                 \frac{J_m(a+ibt)-J_m(a+ib)}{1-t^2}\,dt.
\end{equation}
Then $I_0=S$ and $\partial_a^m I_m=S$ for $m\ge1$.
For a prescribed base point $a_0>0$ the fully anchored primitive is
\begin{align}
 \widehat I_m(a,b;a_0)
 &=I_m(a,b)-\sum_{j=0}^{m-1}
       \frac{(a-a_0)^j}{j!}\partial_a^j I_m(a_0,b)
       \label{gp:eq:anchored}\\
 &=\frac1{(m-1)!}\int_{a_0}^{a}(a-u)^{m-1}S(u,b)\,du.
       \notag
\end{align}
All derivatives of orders below $m$ vanish at $a=a_0$.  The numerator
in \eqref{gp:eq:I} is $O(1-t)$ uniformly on compact parameter sets,
so the displayed integrals and differentiations require no finite-part
interpretation.  Together with \eqref{gp:eq:zeta-second}, this supplies
explicit Hurwitz--Gamma primitives of the nonlocal second-derivative
correction.


% END INCLUDED SOURCE: sections/05_gamma_mass_primitives.tex

% BEGIN INCLUDED SOURCE: sections/06_several_masses.tex
\section{Several quadratic factors: the complete second derivative and confluence}
\label{mm:section}

The seventh further-research question in \cite{QSpec} asks for the
second derivative of a spectrum with several masses and its coincident
limit.  The even-polynomial setting of that report admits a complete
answer.  In contrast to the first derivative, a logarithmic cross term
remains; every coefficient of its coincident-mass expansion nevertheless
has a finite polygamma and Bernoulli evaluation.

\subsection{A finite definition for the weighted logarithmic kernel}

Fix $a>0$ and $P(v)=\sum_{d=0}^{D}p_dv^d$.  For $B,C>0$ set
\begin{equation}\label{mm:eq:beta}
 \beta_k(B,C)=(-1)^k\sum_{j=1}^{k-1}
        \frac{(B^j-C^j)(B^{k-j}-C^{k-j})}{j(k-j)},\qquad k\ge2.
\end{equation}
These are the coefficients in
\[
 \log^2\frac{1+Bz}{1+Cz}=\sum_{k\ge2}\beta_k(B,C)z^k.
\]
They are symmetric in $B,C$ and vanish when $B=C$.  Define the
following specific compensated kernel, with $x_n=n+a$:
\begin{align}
 K_{P,a}(B,C)
 &=\sum_{n\ge0}P(x_n^2)
    \left\{\log^2\frac{x_n^2+B}{x_n^2+C}
                 -\sum_{k=2}^{D}\frac{\beta_k(B,C)}{x_n^{2k}}\right\}
    \notag\\
 &\quad+\sum_{d=0}^{D}\sum_{k=2}^{D}
              p_d\beta_k(B,C)\zeta(2k-2d,a).
              \label{mm:eq:K}
\end{align}
All inner sums are empty when $D<2$.  The infinite sum has
$O(x_n^{-2})$ tails and is ordinary and locally normally convergent.
The finite counterterms have no pole: their Hurwitz orders are even,
hence never one.  For $P=1$ or $P(v)=v$, the unmodified logarithmic
sum already converges absolutely.

\begin{theorem}[All even weights and any number of masses]
\label{mm:thm:second}
For $q\ge1$ and $\mathbf B=(B_1,\ldots,B_q)$ with each $B_j>0$,
let
\begin{align}
 \mathcal Z_{P,\mathbf B}(s;a)
 &=\sum_{n\ge0}P(x_n^2)\prod_{j=1}^{q}(x_n^2+B_j)^{-s},
 \label{mm:eq:product}\\
 Z_{P,j}(s;a)&=\sum_{n\ge0}P(x_n^2)(x_n^2+B_j)^{-s}.
 \label{mm:eq:individual}
\end{align}
Their continuations are holomorphic at zero, and
\begin{equation}\label{mm:eq:second}
 \boxed{\displaystyle
 \mathcal Z_{P,\mathbf B}''(0;a)
 =q\sum_{j=1}^{q}Z_{P,j}''(0;a)
             -\sum_{1\le i<j\le q}K_{P,a}(B_i,B_j).}
\end{equation}
Repeated masses are allowed.  Each coincident pair contributes zero.
\end{theorem}
\begin{proof}
Compare the product spectrum with
$q^{-1}\sum_jZ_{P,j}(qs;a)$.  At each lattice point write
$L_j=\log(x^2+B_j)$.  The constant and first-order coefficients of
\[
 \frac1q\sum_j e^{-qsL_j}-e^{-s\sum_jL_j}
\]
vanish.  Its second derivative at zero is the elementary variance
identity
\begin{equation}\label{mm:eq:variance}
 q\sum_jL_j^2-\left(\sum_jL_j\right)^2
                =\sum_{i<j}(L_i-L_j)^2.
\end{equation}
To justify this calculation after continuation, factor out $x^{-2qs}$
and expand the remaining factors at infinity.  Only even powers of
$x^{-1}$ occur, including after multiplication by $P(x^2)$.  A pole
at $s=0$ would require a resulting $x^{-1}$ term; there is none.
Subtract finitely many asymptotic terms to leave a normally summable
remainder.  The constant and linear coefficients of each difference
term vanish before the Hurwitz factors are differentiated.  Hence its
second derivative is exactly the finite-subtraction continuation of
\eqref{mm:eq:variance}.  The latter is
\eqref{mm:eq:K} for each pair, proving \eqref{mm:eq:second}.
The same finite-subtraction argument proves holomorphy at zero for
the original spectra.  Repeated parameters follow either directly
or from local normal convergence.
\end{proof}

For two unweighted factors, this reads
\[
 \mathcal Z_{1,(B,C)}''(0;a)
 =2Z_{1,B}''(0;a)+2Z_{1,C}''(0;a)
       -\sum_{n\ge0}\log^2\frac{x_n^2+B}{x_n^2+C}.
\]
At $B=C$ it reduces to $4Z_{1,B}''(0;a)$, exactly as required by
$\mathcal Z_{1,(B,B)}(s;a)=Z_{1,B}(2s;a)$.

\subsection{An explicit digamma kernel determines the whole correction}

The elementary quadratic resolvent is
\begin{equation}\label{mm:eq:R}
 R_a(B)=\sum_{n\ge0}\frac1{x_n^2+B}
       =\frac{\psi(a+i\sqrt B)-\psi(a-i\sqrt B)}{2i\sqrt B}.
\end{equation}
For a polynomial weight define its finite continuation
\begin{equation}\label{mm:eq:RP}
 R_{P,a}(B)=P(-B)R_a(B)
    +\sum_{j=1}^{D}p_j\sum_{\ell=0}^{j-1}
                      (-B)^\ell\zeta(-2(j-1-\ell),a).
\end{equation}
It is the continued value $Z_{P,B}(1;a)$.  Formula
\eqref{mm:eq:RP} follows from polynomial division of
$x^{2j}/(x^2+B)$ and the Hurwitz values at nonpositive even integers.
Those values are Bernoulli polynomials:
$\zeta(-m,a)=-B_{m+1}(a)/(m+1)$ \cite{DLMFZeta}.

\begin{theorem}[A normalized double primitive of a digamma divided difference]
\label{mm:thm:integral}
For every $P$ and positive $B,C$,
\begin{equation}\label{mm:eq:double-integral}
 \boxed{\displaystyle
 K_{P,a}(B,C)=\int_C^B\!\int_C^B
              \frac{R_{P,a}(u)-R_{P,a}(v)}{v-u}\,du\,dv.}
\end{equation}
The integrand at $u=v$ is $-R_{P,a}'(u)$.  In particular,
\begin{equation}\label{mm:eq:mixed-derivative}
 \partial_B\partial_CK_{P,a}(B,C)
       =-2\frac{R_{P,a}(B)-R_{P,a}(C)}{C-B}.
\end{equation}
The kernel and both of its first partial derivatives vanish on the
diagonal $B=C$.
\end{theorem}
\begin{proof}
At an ordinary lattice point,
\[
 \log^2\frac{x^2+B}{x^2+C}
   =\int_C^B\!\int_C^B\frac{du\,dv}{(x^2+u)(x^2+v)}.
\]
The rational kernel splits as
\[
 \frac1{(x^2+u)(x^2+v)}
   =\frac{(x^2+u)^{-1}-(x^2+v)^{-1}}{v-u}.
\]
For an arbitrary polynomial weight, perform the same finite
subtractions as in \eqref{mm:eq:K}.  Uniformity on the compact mass
interval permits both integrals to pass through the ordinary remainder;
the finitely many polynomial terms integrate directly.  The resulting
continued rational sum is the divided difference of
\eqref{mm:eq:RP}, proving \eqref{mm:eq:double-integral}.
The diagonal value is removable by differentiation.  Differentiating
the original logarithmic kernel and its finite subtractions gives
\eqref{mm:eq:mixed-derivative}.  The diagonal normalizations follow
from the square or from the double integral.
\end{proof}

\subsection{Every coincident-mass coefficient is finite}

Write $Z_{P}(m;B,a)$ for the continued value of
$\sum_nP(x_n^2)(x_n^2+B)^{-m}$.  It is regular at every positive
integer $m$, since possible spectral poles for an even polynomial
weight occur at half-integers.

\begin{theorem}[The full confluent expansion]
\label{mm:thm:confluence}
For $|h|<a^2+B$,
\begin{equation}\label{mm:eq:confluence}
 \boxed{\displaystyle
 K_{P,a}(B,B+h)=2\sum_{m\ge2}
       \frac{(-1)^mH_{m-1}}m\,h^m Z_{P}(m;B,a).}
\end{equation}
Every coefficient has the finite evaluation
\begin{equation}\label{mm:eq:coefficient-resolvent}
 Z_{P}(m;B,a)=\frac{(-1)^{m-1}}{(m-1)!}
                            \partial_B^{m-1}R_{P,a}(B).
\end{equation}
Consequently no new unspecified constants are introduced in this
coincident-mass expansion: the coefficients belong to the finite
algebra generated by polygamma values at $a\pm i\sqrt B$,
Bernoulli-polynomial values, and rational functions of $B$.
\end{theorem}
\begin{proof}
Use the ordinary Taylor identity, with the logarithm that vanishes at zero,
\[
 \log^2(1+z)=2\sum_{m\ge2}\frac{(-1)^mH_{m-1}}m z^m,
 \qquad |z|<1.
\]
The nearest possible zero of $x_n^2+B+h$ occurs at
$h=-(a^2+B)$.  On every smaller disk the finite-subtracted logarithmic
sum is holomorphic and the same subtractions justify coefficient
extraction.  More explicitly, separate the finitely many indices
$m\le D$, which use continued values, from $m\ge D+1$,
whose weighted sums converge ordinarily and normally on the smaller
disk.  A coefficient is precisely the continued weighted
resolvent $Z_P(m;B,a)$, which proves \eqref{mm:eq:confluence}.
Finally, differentiate the identity
$\partial_B Z_P(s;B,a)=-sZ_P(s+1;B,a)$, first in its convergence
region and then by continuation, starting at $s=1$.  No integer-order
pole is crossed.  This gives \eqref{mm:eq:coefficient-resolvent} and
the asserted finite special-function description.
\end{proof}

Equivalently, polynomial Taylor expansion about $-B$ gives
\[
 Z_P(m;B,a)=\sum_{\ell=0}^{D}\frac{P^{(\ell)}(-B)}{\ell!}
                          Z_1(m-\ell;B,a).
\]
For $j\ge0$ the required nonpositive-order values are simply
\[
 Z_1(-j;B,a)=\sum_{\ell=0}^{j}\binom j\ell
                              B^{j-\ell}\zeta(-2\ell,a),
\]
and the positive orders are derivatives of \eqref{mm:eq:R}.
This second formula is convenient for exact symbolic implementation.

For instance, when $P=1$,
\[
 K_{1,a}(B,B+h)
 =h^2Z_1(2;B,a)-h^3Z_1(3;B,a)
       +\frac{11}{12}h^4Z_1(4;B,a)+O(h^5).
\]
At $a=1/2$, $B=b^2$, the first coefficient is elementary:
\begin{equation}\label{mm:eq:half-explicit}
 Z_1(2;b^2,1/2)
 =\frac{\pi\tanh(\pi b)}{4b^3}
                     -\frac{\pi^2\operatorname{sech}^2(\pi b)}{4b^2}.
\end{equation}
It follows by differentiating
$R_{1/2}(b^2)=\pi\tanh(\pi b)/(2b)$.
The corresponding bilateral results follow by adding the formulas
at $a$ and $1-a$, for $0<a<1$.  In that sum all nonpositive even
Hurwitz values cancel by Bernoulli reflection, leaving an even shorter
weighted resolvent.  The higher spectral derivatives at a generic
noncoincident mass remain described by the exact logarithmic kernels;
no unproved finite polylogarithm reduction is inferred from the finite
confluent coefficients.


% END INCLUDED SOURCE: sections/06_several_masses.tex

% BEGIN INCLUDED SOURCE: sections/07_rank_refinement.tex
% Research fragment. No preamble. All results and scope checked below.
% Primary source inputs: BTD general shifted-Hankel reduction and current
% Polygamma Reconstruction and Resonance, Section 5 (order eight).
\section{Higher precision and the order-ten optimal prefix}
\label{rt:section}

The current reconstruction report proves the order-eight prefix outside
$b\equiv1\pmod{16}$. We retain its exact four-square contraction and
its integral numerator $P_8$; that reduction is inherited from \cite{PRR}. The new
contribution is a higher precision congruence that treats most of that
exceptional class, together with a new five-square contraction and
nonvanishing theorem at order ten. No assertion below settles the
entire gap-six family or all unequal even orders.

\subsection{Notation and the source reduction}

For an integer $s\ge2$, let
\[
 A_s(n)=\zeta(s)-H_n^{(s)}=\zeta(s,n+1),\qquad x_n=n+\tfrac12.
\]
For even orders $p,q$, define
\[
 D_{p,q}(w)=\sum_{n=0}^{\infty}A_p(n)A_q(n)x_n^{-w},
 \qquad \operatorname{Re}w>3-p-q,
 \qquad D_r=D_{p,q}(-2r)
\]
by meromorphic continuation at the last arguments. The source proves
regularity there. Put
\[
 C_j=\sum_{r=0}^j\binom{2j+1}{2r}2^{-2(j-r)}D_r.
\]
Its inverse transform is
\[
 D_r=\sum_{j=0}^r\frac{\binom{2r}{2j}}{2j+1}
                  B_{2r-2j}(1/2)C_j.
\]
For $p=2a$, $q=2b$, $a<b$, the source shows that all late $C_j$,
$j\ge b$, have rational coordinates on
$Z_k=\zeta(2k)$, $0\le k<b$, where $Z_0=\zeta(0)=-1/2$.
Thus the formal coordinate space used here has the $2b$ basis labels
\[
 C_0,\ldots,C_{b-1},\quad Z_0,\ldots,Z_{b-1}.
\]
Their specified numerical values define an evaluation map from this
formal rational vector space; the labels are independent by definition,
and the numerical values are not asserted to be independent.
The exact inverse Bernoulli transforms, late-coordinate interpretation,
and omitted-binomial convention are inherited from \cite{BTDsource,PRR}.
The coordinate matrix of $C_0,\ldots,C_{2b-1}$ is
$\operatorname{diag}(I_b,M_{a,b})$. The displayed inverse transform
from $C$ to $D$ has diagonal entries $1/(2r+1)$, which explains
the full-prefix determinant factor.

Let $p=2a$, $q=2b$, $1\le a<b$, and put
\[
 \mu_j=\begin{cases}B_{2j}/(2j)!,&j\ge0,\\0,&j<0.\end{cases}
\]
We use $B_1=-1/2$; in particular $\mu_0=1$, and the zero values at negative indices are a convention, not analytic continuation.
The source matrix for the late coordinates of the first $q$ centered
moments is
\[
 (M_{a,b})_{tk}=\mu_{t+k-a+1}
 \left\{\frac{(2b+2t-2a+1)!}{(2b-2k)!}
  +\mathbf1_{k<a}\frac{(2t+1)!}{(2a-2k)!}\right\},
 \qquad 0\le t,k<b.
\]
Its full prefix determinant is $\det M_{a,b}/(4b-1)!!$. Let
$H_{a,b}=(\mu_{t+k-a+1})$ and let $U$ be the rank-$a$ perturbation
in the source factorization $M_{a,b}=R(H_{a,b}+U)E$. Define
\[
 \rho_{a,b}=\frac{\det(H_{a,b}+U)}{\det H_{a,b}}.
\]
The source proves $H_{a,b}$ nonsingular \cite{BTDsource}; the diagonal factors $R,E$
are nonzero. Consequently $\rho_{a,b}\ne0$ is equivalent to the
optimal first-$q$ formal-rank assertion. These are rational matrix
identities, not assertions of arithmetic independence of zeta values.

For the finite contraction let $0\le i,\ell<a$ and set
\begin{align*}
 n_i&=b-a+i,&N_{i\ell}&=2(n_i-\ell)+1,&R_\ell&=2b-2a+2+2\ell,\\
 F_{i\ell}&=f_\ell(n_i),&
 f_\ell(n)&=\frac1{(2\ell+1)!}
     \prod_{r=0}^{2\ell-1}(2n-r)(2n+3+r).
\end{align*}
For $W_{i\ell}(z)=(z+2\ell+1)
 (2n_i+2\ell+z+1)^{\underline{2i}}$, put
\[
 d_{i\ell h}=\frac1{h!}\sum_{j=0}^h
            (-1)^{h-j}\binom hj W_{i\ell}(N_{i\ell}-j).
\]
Using independent symbols $U_1,\ldots,U_{2a-1}$ define
\begin{equation}\label{rt:A}
 \mathcal A^{[a]}_{i\ell}
 =2F_{i\ell}+
 \frac{R_\ell!}
 {(2n_i+1)(2n_i+2)(2\ell+2)!N_{i\ell}!}
 \sum_{\substack{1\le h\le2i+1\\h\text{ odd}}}
 d_{i\ell h}N_{i\ell}^{\underline h}
 U_{a-i+\ell+(h-1)/2}.
\end{equation}
Factorial ratios are finite rational functions of $b$, since
$R_\ell-N_{i\ell}=1+4\ell-2i$ is fixed. For $b\ge2a$ the source
small-determinant theorem gives
\begin{equation}\label{rt:ratio}
 \rho_{a,b}=
 \left.\frac{\det\mathcal A^{[a]}}{\det F}
                    \right|_{U_j=2B_{2b-2j}}.
\end{equation}
Here is the normalization check, including the endpoint correction.
The source Pad\'e row transformation, followed by reversal of the
first $a$ columns and division of row $i$ by $c_{n_i}$, takes the
unperturbed entries to $F_{i\ell}$. Expand
$W_{i\ell}(z)=\sum_h d_{i\ell h}(N_{i\ell}-z)^{\underline h}$
in its finite Newton basis and use
\[
 \sum_{j=0}^N\binom Nj B_j(N-j)^{\underline h}
     =N^{\underline h}B_{N-h}(1).
\]
For $b\ge2a$, $N-h\ge2$ for retained odd $h$, and the even $h$
terms have odd index at least three and vanish. Restoring the
omitted $j=1$ term contributes precisely a second $F_{i\ell}$.
This gives \eqref{rt:A}; the factors two in $2F$ and
$U_j=2B_{2b-2j}$ are both necessary. At smaller $b$, the exceptional
$B_0$ and $B_1(1)$ terms require the direct finite calculation.

\subsection{The new five-square integer certificate}

Put
\begin{align*}
 D_{10}(b)={}&4096b^{12}-110592b^{11}+1236992b^{10}-7303680b^9\\
 &+23705088b^8-38140416b^7+14925056b^6+24490080b^5\\
 &-4469984b^4-15838992b^3-2571948b^2+4244400b+1488375.
\end{align*}
For $a=5$, direct finite polynomial identities give
\begin{align}\label{rt:Fdet}
 \det F={}&\frac{16384}{4465125}
 (b-3)(b-2)^2(b-1)(2b-7)(2b-5)^2(2b-3)\notag\\
 &\quad\cdot(4b-15)(4b-13)(4b-11)^2(4b-9)^2(4b-7)^2
 (4b-5)(4b-3)D_{10}(b),\\
 \frac{\det\mathcal A^{[5]}}{\det F}
 ={}&\frac{P_{10}(b;U_1,\ldots,U_9)}
 {138240(2b-7)(2b-5)D_{10}(b)},\label{rt:P10}
\end{align}
where $P_{10}\in\mathbb Z[b,U_1,\ldots,U_9]$. The polynomial has
439 monomials in the nine $U$ variables; their coefficients have
degree at most 24 in $b$. In the full ten-variable expansion there
are 10,403 monomials. Its complete coefficients are supplied with
the exact verification artifact. Equivalently,
\eqref{rt:A}, \eqref{rt:Fdet}, and \eqref{rt:P10} specify it without
interpolation or an omitted fitted formula.

The identities can be verified by clearing the least common denominator
of each row of \eqref{rt:A}, summing the 120 signed permutation
products, and cross-multiplying \eqref{rt:P10}. All coefficient
comparisons are over the rational polynomial ring. The delivered
verifier performs these operations exactly and checks the integer
coefficients of $P_{10}$. Since $D_{10}(b)$ is odd at every integer
$b$, the denominator in \eqref{rt:P10} always has valuation 10.

For order eight we retain the current source polynomial \cite{PRR}, normalized by
\[
 \rho_{4,b}=
 \frac{P_8(b;2B_{2b-2},\ldots,2B_{2b-14})}
 {576(2b-5)D_8(b)},\qquad b\ge8,
\]
where
\[
 D_8(b)=256b^8-3584b^7+18496b^6-40544b^5+25984b^4
       +18424b^3-5856b^2-8316b-2025.
\]
Its denominator has valuation 6. The independently regenerated
$P_8$ agrees coefficientwise with the source certificate.

\subsection{A uniform Bernoulli congruence at prescribed precision}

All congruences are in $\mathbb Z_{(2)}$, the rationals with odd
denominator. We use the source's elementary fact
$v_2(B_{2j})=-1$ for $j\ge1$; all Bernoulli numbers have valuation
at least $-1$.

\begin{lemma}\label{rt:general-B}
Let $K\ge4$ and $m\ge1$ be integers, let $s=\lceil K/2\rceil$, and suppose
$2m\ge K+s-1$. Set $J=\lfloor(K+s-2)/3\rfloor$. Then
\[
 2B_{2m}\equiv \frac1{2^{s-1}}
 \sum_{\substack{1\le a<2^s\\a\text{ odd}}}
 \sum_{j=0}^{J}\binom mj(a^2-1)^j\pmod{2^K}.
\]
The right side is a finite integer-valued polynomial in $m$.
\end{lemma}
\begin{proof}
Faulhaber's identity gives
\[
 \sum_{a=0}^{2^s-1}a^{2m}
 =2^sB_{2m}+
 \sum_{h=2}^{2m+1}\binom{2m}{h-1}\frac{2^{sh}}h
                         B_{2m+1-h}.
\]
The $h=2$ term is zero. For $h\ge3$, its valuation is at least
$sh-v_2(h)-1\ge s+K-1$. Thus division by $2^{s-1}$ proves
$2B_{2m}\equiv2^{1-s}\sum a^{2m}\pmod{2^K}$. The even summands
vanish at this precision because $2m\ge K+s-1$.
For odd $a$, $a^2-1$ is divisible by eight. Expanding
$a^{2m}=(1+(a^2-1))^m$ and discarding $j>J$ gives an error of
valuation at least $3j-s+1\ge K$ in each term. This proves the
formula. For completeness, sums of even powers of the odd integers below $2^s$ are divisible by $2^{s-1}$: pair $a$ with $a+2^{s-1}$ and induct on $s$. Thus the coefficients in the binomial basis are integers. The finite average may also be retained in its displayed
rational form; no claim about dividing a congruence by a nonunit
without increasing its precision is being made.
\end{proof}

For $K=16$, $s=8$, finite coefficient simplification gives the
particularly useful integer polynomial
\begin{equation}\label{rt:B16}
 \begin{split}
 \mathcal B(x)={}&1+20068x-26384x^2+2816x^3-19712x^4
                  +4096x^5-24576x^6,\\
 2B_{2m}\equiv{}&\mathcal B(m)\pmod{65536},\qquad m\ge12.
 \end{split}
\end{equation}
For an entirely finite specification of this simplification, expand
\[
 \frac1{128}\sum_{a=1,3,\ldots,255}
       \sum_{j=0}^{7}\binom{x}{j}(a^2-1)^j.
\]
Every coefficient of its difference from $\mathcal B(x)$ has
valuation at least 16. The exact program checks this in
$\mathbb Q[x]$, including the denominators from $\binom{x}{j}$.
The degree-seven coefficient vanishes at the displayed precision.

\subsection{Two short congruences and their unique exceptional classes}

Define
\begin{align*}
 R_8(b)={}&4096b^9-5120b^8+4864b^7+2176b^6-176b^5\\
 &-5184b^4-4912b^3-7948b^2+1195b-1727,\\
 R_{10}(b)={}&-512b^8+128b^7-928b^6-816b^5+464b^4
                  +352b^3+118b^2+217b+694.
\end{align*}
The following are coefficientwise polynomial congruences:
\begin{align}
 P_8\bigl(b;\mathcal B(b-1),\ldots,\mathcal B(b-7)\bigr)
       &\equiv4R_8(b)\pmod{65536},\label{rt:R8}\\
 P_{10}\bigl(b;\mathcal B(b-1),\ldots,\mathcal B(b-9)\bigr)
       &\equiv32R_{10}(b)\pmod{65536}.\label{rt:R10}
\end{align}
They follow by direct integer polynomial substitution in the finite
certificates; the accompanying program performs exact modular
polynomial multiplication and verifies every coefficient. After the
actual Bernoulli coordinates are substituted, the precision after division
is explicitly
\begin{align*}
 \frac{P_8(b;2B_{2b-2},\ldots,2B_{2b-14})}{4}
       &\equiv R_8(b)\pmod{16384} &&(b\ge19),\\
 \frac{P_{10}(b;2B_{2b-2},\ldots,2B_{2b-18})}{32}
       &\equiv R_{10}(b)\pmod{2048} &&(b\ge21).
\end{align*}

All coefficients of degree at least two in either $R_8$ or $R_{10}$
are even, and the linear coefficients are odd. Consequently, for
all integers $x,y$,
\begin{equation}\label{rt:isometry}
 v_2(R_j(x)-R_j(y))=v_2(x-y),\qquad x\ne y,\quad j=8,10.
\end{equation}
Indeed, after factoring $x-y$, the remaining polynomial is the odd
linear coefficient plus an even integer. In particular, a root
modulo $2^L$ is unique if it exists. Exact substitution gives
\[
 R_8(8001)\equiv0\pmod{16384},\qquad
 R_{10}(962)\equiv0\pmod{2048}.
\]
Thus these are the unique exceptional classes at their respective
precisions. Their existence and uniqueness do not rely on testing
all residues.

\begin{theorem}[Order-eight refinement]\label{rt:rank8}
For every integer $b\ge5$ with $b\not\equiv8001\pmod{16384}$,
$M_{4,b}$ is nonsingular. Equivalently, the first $q=2b$ centered
moments have full formal rank $q$ for $(p,q)=(8,2b)$. For $b\ge19$
in this range,
\[
             v_2(\rho_{4,b})=v_2(b-8001)-4.
\]
\end{theorem}
\begin{proof}
For $b\ge19$, each $b-j\ge12$ in the seven Bernoulli coordinates.
Apply \eqref{rt:B16} and \eqref{rt:R8}. If $b$ is not in the
exceptional class, \eqref{rt:isometry} gives
$v_2(R_8(b))=v_2(b-8001)<14$. Hence the true numerator has
valuation $2+v_2(b-8001)$, and division by the denominator of
valuation 6 proves the assertion. The existing order-eight theorem
handles $5\le b\le18$ except $b=17$. Direct calculation gives
\[
 \rho_{4,17}=-\frac{5564518155020169678178976040991}
                         {8937041637533625}\ne0.
\]
\end{proof}

\begin{theorem}[Order-ten prefix]\label{rt:rank10}
For every integer $b\ge6$ with $b\not\equiv962\pmod{2048}$,
$M_{5,b}$ is nonsingular. Equivalently, the first $q=2b$ centered
moments have full formal rank $q$ for $(p,q)=(10,2b)$. For $b\ge21$
in this range,
\[
              v_2(\rho_{5,b})=v_2(b-962)-5.
\]
\end{theorem}
\begin{proof}
For $b\ge21$, all nine indices $b-j$ are at least 12. Equations
\eqref{rt:B16}, \eqref{rt:R10}, and \eqref{rt:isometry} give
$v_2(P_{10})=5+v_2(b-962)$ whenever $b\not\equiv962\pmod{2048}$.
The denominator in \eqref{rt:P10} has valuation 10. This proves
nonvanishing and the valuation formula.
For $6\le b\le20$, an exact finite calculation with the displayed
source matrix gives the respective valuations
\[
 -4,-4,-6,-5,-2,-5,-4,-5,-3,-5,-4,-5,-1,-5,-4.
\]
The complete nonzero determinants and ratios are recorded in
\texttt{direct\_rank\_checks.json}; its independent verifier constructs
$M_{5,b}$ and $H_{5,b}$ from Bernoulli numbers and takes their exact
rational determinants. It does not invoke the generic contraction
below its valid threshold $b\ge10$.
\end{proof}

For example, the previously unsettled gap-six pair $(p,q)=(10,16)$
is included among these finite base cases, with
\[
 \det M_{5,8}=\frac{537653622905129}{68669669376000}>0.
\]
The first pair in the excluded order-ten progression has $b=962$;
its exclusion is a limitation of the congruence, not evidence of a
zero determinant. The same applies to $b=8001$ in the order-eight
refinement. The global gap-six conjecture remains open here.

\subsection{A useful lower precision corollary and exact scope}

The same finite calculation at modulus $512$ gives
\[
 P_{10}\bigl(b;\mathcal B_9(b-j)\bigr)
 \equiv32(6b^2-7b+6)\pmod{512},
\]
where
$\mathcal B_9(x)=256x^4+256x^3+240x^2+100x+1$
represents $2B_{2m}$ modulo $512$ for $m\ge7$. The odd linear
coefficient again gives a unique zero, now $b\equiv2\pmod{16}$.
Thus a shorter statement is that order ten is proved outside
$b\equiv2\pmod{16}$, with
$v_2(\rho_{5,b})=v_2(b-2)-5$ for $b\ge16$ outside that class.
The higher precision theorem treats many additional values within it.

Further increases in precision are algorithmic through
Lemma~\ref{rt:general-B}, but an arbitrarily long list of certified
residue classes is not a proof for all $b$. In particular, neither
the order-eight nor the order-ten formula above asserts a global
valuation law obtained by replacing the exceptional class with a
single ordinary integer. A proof for the remaining progressions,
or an independent archimedean nonvanishing argument, is still needed.

% END INCLUDED SOURCE: sections/07_rank_refinement.tex

% BEGIN INCLUDED SOURCE: sections/08_audit_and_research.tex
\section{Source reconciliation and further research}
\label{audit:section}

\subsection{Corrections, qualifications, and status changes}

\paragraph{Current source status.}
The targeted audit did not establish an additional false theorem in the
current source dependencies. The useful corrections are therefore
precise integration and scope corrections, not a claim to have found
a new error merely because a stronger statement is unavailable.
The two Stieltjes corrections already incorporated in the unified
manuscript must be retained. Its chapter
\path{11-stieltjes-source-corrections.tex} explicitly repairs an
initial-index discrepancy in the inspected version of a harmonic-zeta
paper and a literal-coefficient sign description in a height-one
multiple-zeta paper \cite{ProveIt}. Neither is a new discovery of this
delivery.

The exact conventions can be checked without a numerical fit. From
\eqref{intro:stieltjes} and the Hurwitz recurrence,
\[
 \gamma_r(a+1)=\gamma_r(a)-\frac{\log^r a}{a}.
\]
Consequently a Stieltjes limit whose sum begins at $n=1$ instead of
$n=0$ has lost $\log^r(a)/a$. For $r=0,a=1$ the lost quantity is
one. Likewise the literal coefficient of $u$ in
$\zeta(1+u,a)$ is $-\gamma_1(a)$, rather than $\gamma_1(a)$.
These facts explain both inherited corrections and fix the sign
convention needed for \eqref{ats:eq:stieltjes} and
\eqref{sd:eq:resonant}.

\paragraph{Higher-order pole terms.}
For a product $C_k(s)\zeta(1+2s,a)$, the coefficient of $s^r$
uses the coefficient of $s^{r+1}$ in $C_k$. Formula
\eqref{sd:eq:resonant} retains this term as
$r!\alpha_{k,r+1}/2$. A proposed extension obtained by differentiating
only the regular Laurent part would be incorrect. The present article
supplies the complete replacement, without suggesting that the
source's proved first-derivative theorem was wrong.

\paragraph{Ordinary sums and continued values.}
The continuous-order square law is an ordinary identity only on its
stated half-plane. At $s=0$ and a fixed positive $a$, continuation gives
$\mathcal B_0(z)=1/2$. The continued lower-end expression in
\eqref{ats:eq:continuous} is then $1/4$, although every pointwise
summand of its left side at $s=0$ is zero. In the finite identity
\eqref{ats:eq:continuous-finite}, the upper endpoint is exactly $1/4$
and cancels it. This is a concrete obstruction to silently extending
the ordinary summation prescription across that boundary. Similarly,
the finitely many low-order coefficients in
\eqref{mm:eq:confluence} are continued spectral values when the
unmodified weighted series diverges. The article states that distinction
explicitly.

\paragraph{Moment reconstruction.}
The gap-four theorem is already proved in \cite{HTGC,HTST}. An older
research-status paragraph that still lists it as wholly open should
be updated by citing those reports. For order eight, the inherited
theorem outside $b\equiv1\pmod{16}$ is strengthened here; it is not
relabelled as an original result. The new statements are the smaller
exceptional class in Theorem~\ref{rt:rank8} and the order-ten family in
Theorem~\ref{rt:rank10}.

Experimental valuation laws must be replaced by their proved
finite-precision versions. For example, the new order-eight formula
gives valuations $4$ at $b=65$ and $2$ at $b=129$, while the
order-ten formula gives valuation $2$ at $b=66$.
These follow from $b-8001$ and $b-962$, respectively. They do not
justify a global law centered at an ordinary integer: the excluded
progressions remain outside the theorem. Moreover, formal full rank
refers to the specified moment/zeta coordinate module. It must not
be converted into arithmetic independence of its numerical evaluations.

\paragraph{Cyclotomic parity.}
The finite coefficient identity in Section~\ref{cyc:sec} is a form
of the classical colored depth-two parity theorem \cite{PanzerParity}.
The addition is its explicit rational-tail resummation with arbitrary
outer character, its cutoff proof, and the specialized evaluations.
The complementary conjugation sector is not reduced. Conditional
bilateral intermediate sums are locally uniformly convergent; the
projected double sums used in the final generator are absolutely
normally convergent. Replacing the first statement by absolute normal
convergence would be an unjustified strengthening.

\paragraph{Gaussian conjectures and formal obstructions.}
The source's Gaussian weight-six and revised weight-eight conjectures
remain unresolved here. A non-membership result in a specified
shuffle, parity, or Cayley relation space is a limitation of that
relation system. It is not a disproof of an equality between actual
periods. The newly evaluated Gaussian rational tails have their own
definitions and cannot be substituted for those older conjectural
quantities. No integration note should mark the older problems solved
on the strength of the present conductor-four examples.

\subsection{A proposed research programme with concrete endpoints}

Each question below asks for an identity, a complete structural
description, or an exact nonvanishing certificate. The order reflects
dependence on the proved results, not a claim about difficulty.

\begin{question}[The two remaining moment progressions]
Can one prove nonvanishing in the excluded classes
$b\equiv8001\pmod{16384}$ for order eight and
$b\equiv962\pmod{2048}$ for order ten?
Lemma~\ref{rt:general-B} makes higher congruence precision effective.
A useful next result would either control the exceptional classes
uniformly at every precision or give an independent exact argument on
them. Certifying any fixed number of additional binary digits still
leaves an infinite progression. One must also distinguish a compatible
$2$-adic root from an ordinary positive integer at which the rational
determinant actually vanishes.
\end{question}

\begin{question}[The full gap-six family]
For every even $p=2a$ and $q=2a+6$, does the first-$q$ prefix have
full formal rank? The finite example $(p,q)=(10,16)$ is included here,
but a proof uniform in $a$ is still missing. The fixed-order contraction
has the threshold $b\ge2a$, whereas here $b=a+3$. For $a>3$ its
omitted $B_0$ and $B_1$ contributions cannot be discarded. An exact
bordered-Hankel formula that retains those moving endpoint terms is
a concrete next target.
\end{question}

\begin{question}[All unequal even orders]
Find a uniform nonvanishing theorem for $M_{a,b}$, $1\le a<b$,
or identify and prove the exact exceptional set if one exists.
A useful intermediate result would characterize the small determinant
after the Pad\'e transformation in terms of explicit orthogonal
polynomial evaluations and their finite differences. The endpoint
should be an exact identity or sign/nonzero certificate, rather than
an extrapolation from a large table of ranks.
\end{question}

\begin{question}[Bilinear adjacent-tail certificates]
Replace $E_n^2$ in Theorem~\ref{ats:thm:polynomial} by a product of
tails with distinct polynomial denominators $P,Q$. Allow both products
of adjacent tails in the certificate. Which rational coefficient
functions eliminate every linear-tail term, and which pairs $P,Q$
admit a finite endpoint evaluation? The compatibility calculation
in Section~\ref{ats:section} classifies only the displayed one-product
ansatz. The first nontrivial test is two distinct quadratic
resolvents with independent shifts.
\end{question}

\begin{question}[Continuation of the ordinary square law]
Continue \eqref{ats:eq:continuous} across $\Re s=0$ with a fully
specified subtraction of its exact upper endpoint. Determine the
contact term and all order derivatives at each nonpositive integer.
The $s=0$ discrepancy of $1/4$ described above is a mandatory first
test. An all-index formula would connect negative-order alternating
Hurwitz values, Euler polynomials, and Laurent coefficients without
confusing ordinary sums with continued values.
\end{question}

\begin{question}[Independent shifts and independent order jets]
The two factors of $P(k)=(k+a)(k+b)$ in
\eqref{ats:eq:two-shift} belong to one tail. They do not answer the
source question about
$\zeta(p,n+v)\zeta(q,n+w)$ with independently shifted tails.
Derive a finite-cutoff identity for those products, then differentiate
independently in $p,q$, and the outer weight order. The desired formula
must state all moving polar hyperplanes and their residues. This would
extend the matched continuous-order diagonal treated here and resolve
the corresponding independent-shift question in \cite{HTST}.
\end{question}

\begin{question}[The complementary cyclotomic sector]
Construct the two-resolvent generator for
\[
 \frac{R_{a,b;c}(x,y)-R_{a^{-1},b^{-1};c^{-1}}(x,y)}{2i}
\]
in the odd-weight setting. A first conductor-four coefficient is
$\Im Z(2,1;i,-i)$. Specify a candidate collection of classical or
cyclotomic polylogarithms and an exact reduction into it, or retain a
precisely defined additional coordinate. The parity generator alone
does not determine this sector. Failure of a finite numerical search
would not prove independence from the proposed collection.
\end{question}

\begin{question}[Shifted and degree-generating cyclotomic normal forms]
Introduce independent rational shifts into the bilateral lattice proof
and compute its finite boundary terms. Then introduce a third generating
variable that sums the outer polynomial degree $h$ in
\eqref{cyc:allweights}. Having a formula for every individual $h$ does
not yet produce that analytic generating function. A compact answer
would organize all resolvent and weight derivatives by a finite
Bell-polynomial rule, with its convergence region stated.
\end{question}

\begin{question}[Three-tail products]
Develop a three-tail analogue of the finite outer-cutoff decomposition.
The natural target is an explicit depth-two remainder plus finite
single-polylogarithm terms and endpoints. Determine the exact parity
sector in which this reduction occurs. The product of three tails
has a new ordered region structure, so unrestricted depth-one closure
would be an unsupported extrapolation from Section~\ref{cyc:sec}.
\end{question}

\begin{question}[A finite description of the logarithmic Gamma correction]
At which rational shifts and nonzero masses can $S(a,b)$ in
\eqref{gp:eq:S} be expressed in a stated finite collection of
polylogarithm values and order derivatives? The elementary-kernel
forms \eqref{gp:eq:integer}--\eqref{gp:eq:half} and the exact
shift/distribution rules provide starting equations. A rigorous
finite evaluation at a single nontrivial mass would be progress.
For generic parameters, the Gamma integral remains a complete exact
representation; no minimality theorem for that representation has
been established.
\end{question}

\begin{question}[Several masses at every spectral order]
Extend Theorem~\ref{mm:thm:second} to $r\ge3$. Pointwise, the variance
identity is replaced by a higher symmetric polynomial in the logarithms
$L_j$. Determine the finite-subtracted kernels, their confluent
coefficients, and their mixed mass primitives. With even weights
there is no pole at zero; with odd weights, compute the additional
Stieltjes and one-higher-coefficient residue terms rather than assuming
the even-weight prescription still applies.
\end{question}

\begin{question}[Shift, mass, and order derivatives in one formula]
Combine the primitive tower $J_m$ with the divided-difference kernel
$R_{P,a}$. Can every rectangular mixed shift--mass primitive of the
second defect be written using a finite, explicitly anchored collection
of Hurwitz derivatives and ordinary integrals? The mixed derivatives
must commute on a stated domain, and constants must be fixed at all
base faces. This is a stronger target than the separate one-variable
and double-mass primitives proved here.
\end{question}

\begin{question}[Joint tail identities and spectral identities]
The square certificate, the cyclotomic parity generator, and the
quadratic resolvent all produce polygamma expressions. At a fixed
conductor and weight, form their joint exact relation system and
determine which mixed harmonic-tail coefficients it isolates. A
particularly concrete endpoint is a new closed formula for a single
mixed tail product, rather than only the convolution in
\eqref{ats:eq:convolution}. Keep the formal relation quotient distinct
from the space of numerical periods.
\end{question}

\begin{question}[Small trusted certificates and formalization]
Which of the finite proofs can be reduced to short rational identities
that a proof assistant can check directly? The adjacent-tail recurrence,
the harmonic coefficient formula, the resonant Laurent multiplication,
and the two modular substitution certificates are natural first
targets. The analytic continuation lemmas then need explicit locally
uniform majorants and an agreed special-function library. A formal
check of the finite data alone would not formalize the associated
infinite-dimensional theorem.
\end{question}


% END INCLUDED SOURCE: sections/08_audit_and_research.tex

\clearpage
% BEGIN INCLUDED SOURCE: sections/09_verification.tex
\section{Reproducibility and integration}
\label{verify:section}

\subsection{Exact algebra and independent numerical diagnostics}

The companion programs use Python, SymPy, and mpmath. They are
self-contained relative to the delivered certificate files and require
no private paths or network access. The recorded environment uses
Python 3.12.14, SymPy 1.14.0, and mpmath 1.3.0. The proofs themselves
do not depend on a particular floating-point library.

\begin{longtable}{@{}>{\raggedright\arraybackslash}p{0.27\textwidth}>{\raggedright\arraybackslash}p{0.36\textwidth}>{\raggedright\arraybackslash}p{0.31\textwidth}@{}}
\toprule
Component & Exact verification & Numerical diagnostics\\
\midrule
\endfirsthead
\toprule Component & Exact verification & Numerical diagnostics\\\midrule
\endhead
Tail square and Stieltjes identities & 96 exact assertions, including
the finite endpoint and deliberately incorrect color-sign rejection &
29 diagnostics at 75 decimal digits; independent integral, finite-cutoff,
and order-derivative calculations.\\[5pt]
Cyclotomic generator & Written bilateral proof and finite parity
coefficient extraction in Section~\ref{cyc:sec} & 12 diagnostics at
60--65 digits, including trivial colors, arbitrary outer character,
odd conductor, rational examples, and the Gamma primitive.\\[5pt]
Spectral defects and mass kernels & 194 exact assertions for harmonic
coefficients, Laurent multiplication, logarithmic kernels, variance,
and confluence & 21 diagnostics at 65 digits, including independent
spectral continuation, Gamma integrals, mass differentiation, and
anchored primitives.\\[5pt]
Moment-rank refinements & Complete rational polynomial determinant
replay, coefficientwise modular substitutions, and original finite
source-matrix determinants & No floating rank test is used as evidence
for a nonvanishing theorem.\\
\bottomrule
\end{longtable}

The 62 high-precision checks complement the 290 exact assertions in
the two analytic verification suites. The larger determinant
certificates are a separate exact calculation, not included in that
assertion count. The order-ten numerator has 10,403 integer monomials
in its full expansion. The verifier reconstructs the determinant from
the displayed rational matrix and cross-multiplies coefficientwise;
it does not infer a polynomial by interpolation at sample values.

The numerical spectral comparison does not merely evaluate the same
closed form twice. It constructs the quadratic zeta by a separate
Pochhammer--Hurwitz expansion, differentiates the finite linear-factor
Hurwitz expression, and compares their difference with the
finite-subtracted logarithmic formula. The tail checks use finite
endpoint identities and independent integral representations where
appropriate. Nonetheless, all decimal results remain diagnostics.
The recorded error tolerances are not rigorous enclosures of the
infinite sums.

The determinant proof was independently reviewed against the original
binomial source matrix, in addition to the factorial form used in the
delivered verifier. The Bernoulli thresholds, division of moduli,
source normalization, and exceptional classes were also checked
independently. The exact replay and the written proof have different
roles: the replay validates finite algebra; the uniform congruence
argument makes that finite algebra apply to infinitely many $b$.

\subsection{Files and commands}

The package has two equivalent \LaTeX\ entry points:
\path{article.tex} inputs the modular sections, while
\path{Alternating_Tails_and_Spectral_Defects.tex} contains the complete
standalone article. The supplied PDF is compiled from the standalone
source. The modular files make thematic integration straightforward.
The bibliography is local to the package; no external bibliography
database is required.

From the extracted package directory, the verification entry point is
\begin{verbatim}
python verification/run_all.py
\end{verbatim}
The exact determinant replay may take longer than the small symbolic
checks. A diagnostic shortcut, which intentionally omits reconstruction
of the two multivariate determinant identities, is
\begin{verbatim}
python verification/verify_rank_refinement.py --quick
\end{verbatim}
Quick mode is useful during editing; the recorded delivery was checked
with the full replay. The programs rewrite their result JSON files.
For an independent review that preserves the delivered evidence,
run them in a copy of the extracted package.

The PDF can be rebuilt with
\begin{verbatim}
latexmk -pdf -interaction=nonstopmode -halt-on-error \
    Alternating_Tails_and_Spectral_Defects.tex
\end{verbatim}
The package includes a short build helper that first regenerates the
standalone source from the modular files. The only nonstandard Python
dependencies are pinned in \path{requirements.txt}. No source
repository checkout is required to reproduce the new certificates.

\subsection{Suggested placement and claim ledger}

This continuation belongs with the harmonic-tail and spectral/Stieltjes
thematic reports under
\path{Analysis/Polylogarithms/docs/reports/}. The delivered directory
can be preserved intact and cross-referenced from both themes. Later
editorial integration into the unified manuscript can use
Sections~\ref{ats:section}--\ref{cyc:sec} for tail identities,
Sections~\ref{sd:section}--\ref{mm:section} for spectral derivatives
and primitives, and Section~\ref{rt:section} for moment rank.
The current incoming procedure treats placement and manuscript
rewriting as separate stages \cite{Incoming}.

The machine-readable claim ledger records what is proved, inherited,
partially advanced, and still open. In particular it keeps the
Gaussian conjectures, complementary cyclotomic sector, independent
tail shifts, global gap-six family, and two exceptional moment
progressions open. The provenance record names the exact repository
pin and the six inspected archive sources. The integration notes
identify the source research questions updated by the new theorems
and preserve the earlier correction history.


% END INCLUDED SOURCE: sections/09_verification.tex

\clearpage
% BEGIN INCLUDED SOURCE: references.tex
\begin{thebibliography}{99}
\addcontentsline{toc}{section}{References}

\bibitem{ProveIt}
V. Reshetnikov, \emph{ProveIt: Polylogarithms unified manuscript},
repository snapshot \shortpin, 11 October 2026.
\href{https://github.com/VladimirReshetnikov/ProveIt/tree/73f18ef013d5b74eafc883efb2c15210642df12a/Analysis/Polylogarithms/docs/manuscript}{Pinned manuscript directory}.

\bibitem{Incoming}
\emph{ProveIt incoming reports}, six archives at repository snapshot
\shortpin, 11 October 2026.
\href{https://github.com/VladimirReshetnikov/ProveIt/tree/73f18ef013d5b74eafc883efb2c15210642df12a/docs/incoming}{Pinned incoming directory}.

\bibitem{HTST}
\emph{Harmonic Tails, Shifts, and Twists}, research continuation for the
ProveIt programme, 11 October 2026. Delivered source
\path{ProveIt_Harmonic_Tails_Shifts_and_Twists_2026-10-11.zip},
in \cite{Incoming}; especially the twisted-tail generator and Q6.

\bibitem{HTGC}
\emph{Harmonic-Tail and Gamma Continuation}, research continuation for
the ProveIt programme, 11 October 2026. Delivered source
\path{ProveIt_Harmonic_Tail_Gamma_Continuation_2026-10-11.zip},
in \cite{Incoming}; gap-four moment-prefix theorem and source status.

\bibitem{QSpec}
\emph{Quadratic-Spectrum Stieltjes Identities: Complete reflected jets,
harmonic counterterms, finite polylogarithmic reductions, and a
polynomial parity criterion}, research continuation for the ProveIt
programme, 10 October 2026. Delivered source
\path{ProveIt_Quadratic_Spectrum_Stieltjes_2026-10-10.zip},
in \cite{Incoming}; further-research questions 2 and 7.

\bibitem{PRR}
\emph{Polygamma Reconstruction and Resonance}, research continuation for
the ProveIt programme, 11 October 2026. Delivered source
\path{ProveIt_Polygamma_Reconstruction_and_Resonance_2026-10-11.zip},
in \cite{Incoming}; Section 5 and its exact order-eight certificate.

\bibitem{BTDsource}
\emph{Bernoulli, Tornheim, and digamma research}, ProveIt research
continuation, 11 October 2026. General shifted-Hankel reduction and
finite bordered determinant, especially
\path{sections/03b_general_hankel.tex} in the delivered archive.
\href{https://github.com/VladimirReshetnikov/ProveIt/blob/446dca9bd65c20e409d1d2561d4ffad408103e2a/docs/incoming/ProveIt_Bernoulli_Tornheim_Digamma_2026-10-11.zip}{Archive at its inspected earlier revision}.
The reduction is also used in the current reconstruction report
\cite{PRR} and the moment chapters of \cite{ProveIt}.

\bibitem{Hoffman}
M. E. Hoffman, \emph{Sums of products of Riemann zeta tails},
arXiv:1602.03399 (2016).
\url{https://arxiv.org/abs/1602.03399}.

\bibitem{PanzerParity}
E. Panzer, \emph{The parity theorem for multiple polylogarithms},
arXiv:1512.04482, version 3.
\url{https://arxiv.org/abs/1512.04482}.

\bibitem{NielsenBeta}
C. Berg, S. Koumandos, and H. L. Pedersen,
\emph{Nielsen's beta function and some infinitely divisible
distributions}, arXiv:1905.04131 (2019).
\url{https://arxiv.org/abs/1905.04131}.

\bibitem{Dowker}
J. S. Dowker, \emph{Calculation of the multiplicative anomaly},
arXiv:1412.0549, version 3.
\url{https://arxiv.org/abs/1412.0549}.

\bibitem{DLMFZeta}
NIST Digital Library of Mathematical Functions, Section 25.11,
\emph{Hurwitz zeta function}; definitions, continuation, parameter
derivatives, and special values. Accessed 11 October 2026.
\url{https://dlmf.nist.gov/25.11}.

\bibitem{DLMFPolylog}
NIST Digital Library of Mathematical Functions, Section 25.12,
\emph{Polylogarithms}. Accessed 11 October 2026.
\url{https://dlmf.nist.gov/25.12}.

\bibitem{DLMFPolygamma}
NIST Digital Library of Mathematical Functions, Section 5.15,
\emph{Polygamma functions}. Accessed 11 October 2026.
\url{https://dlmf.nist.gov/5.15}.

\bibitem{DLMFGamma}
NIST Digital Library of Mathematical Functions, Sections 5.4--5.5,
\emph{Gamma function: special values and functional relations};
in particular Gauss's rational digamma formula (5.4.19), reflection,
and multiplication. Accessed 11 October 2026.
\url{https://dlmf.nist.gov/5.4}; \url{https://dlmf.nist.gov/5.5}.

\end{thebibliography}


% END INCLUDED SOURCE: references.tex

\end{document}

